\documentclass[11pt]{article}

\usepackage[letterpaper,margin=0.82in,headheight=15pt]{geometry}
\usepackage[T1]{fontenc}
\usepackage{lmodern}
\usepackage{microtype}
\usepackage{amsmath,amssymb,amsthm,bm,mathtools}
\usepackage{booktabs,longtable,tabularx,array,multirow}
\usepackage{enumitem}
\usepackage{xcolor}
\usepackage{fancyhdr}
\usepackage{graphicx}
\usepackage{tikz-cd}
\usepackage{hyperref}
\usepackage[nameinlink,noabbrev]{cleveref}

\definecolor{navy}{HTML}{17324D}
\definecolor{teal}{HTML}{147D92}
\definecolor{orange}{HTML}{D47A27}
\definecolor{softblue}{HTML}{EEF4F7}
\definecolor{softgray}{HTML}{F3F5F6}
\definecolor{darkgray}{HTML}{46545E}
\definecolor{softorange}{HTML}{FFF4E8}

\hypersetup{
  colorlinks=true,
  linkcolor=navy,
  citecolor=teal,
  urlcolor=teal,
  pdftitle={Volume III: Dynamical Wilson Lines, Rescattering Phases, and Naive-T-Odd GTMDs},
  pdfauthor={DeuteronWigner project}
}

\pagestyle{fancy}
\fancyhf{}
\fancyhead[L]{\small\color{darkgray}DeuteronWigner formalism - Volume III}
\fancyhead[R]{\small\color{darkgray}Dynamical Wilson lines and phases}
\fancyfoot[L]{\small\color{darkgray}29 July 2026}
\fancyfoot[R]{\small\color{darkgray}\thepage}
\renewcommand{\headrulewidth}{0.4pt}

\setlist[itemize]{leftmargin=1.5em,itemsep=2pt,topsep=3pt}
\setlist[enumerate]{leftmargin=1.7em,itemsep=2pt,topsep=3pt}
\setlength{\parindent}{0pt}
\setlength{\parskip}{5pt}
\setlength{\emergencystretch}{2em}
\renewcommand{\arraystretch}{1.16}
\allowdisplaybreaks

\newtheorem{definition}{Definition}[section]
\newtheorem{axiom}[definition]{Axiom}
\newtheorem{principle}[definition]{Design principle}
\newtheorem{requirement}[definition]{Requirement}
\newtheorem{proposition}[definition]{Proposition}
\newtheorem{theorem}[definition]{Theorem}
\newtheorem{lemma}[definition]{Lemma}
\newtheorem{corollary}[definition]{Corollary}
\newtheorem{remark}[definition]{Remark}

\crefname{definition}{definition}{definitions}
\crefname{axiom}{axiom}{axioms}
\crefname{principle}{principle}{principles}
\crefname{requirement}{requirement}{requirements}
\crefname{proposition}{proposition}{propositions}
\crefname{theorem}{theorem}{theorems}
\crefname{lemma}{lemma}{lemmas}
\crefname{corollary}{corollary}{corollaries}
\crefname{remark}{remark}{remarks}

\newcommand{\kT}{\bm{k}_{T}}
\newcommand{\ellT}{\bm{\ell}_{T}}
\newcommand{\pT}{\bm{p}_{T}}
\newcommand{\qT}{\bm{q}_{T}}
\newcommand{\PT}{\bm{P}_{T}}
\newcommand{\DT}{\bm{\Delta}_{T}}
\newcommand{\DNT}{\bm{\Delta}_{N T}}
\newcommand{\bD}{\bm{b}_{\Delta}}
\newcommand{\bM}{\bm{b}_{\mathrm{TMD}}}
\newcommand{\RT}{\bm{R}_{T}}
\newcommand{\dd}{\mathrm{d}}
\newcommand{\Tr}{\operatorname{Tr}}
\newcommand{\tr}{\operatorname{tr}}
\newcommand{\Span}{\operatorname{span}}
\newcommand{\End}{\operatorname{End}}
\newcommand{\Hom}{\operatorname{Hom}}
\newcommand{\Id}{\operatorname{Id}}
\newcommand{\Ker}{\operatorname{Ker}}
\newcommand{\Ran}{\operatorname{Ran}}
\newcommand{\rank}{\operatorname{rank}}
\newcommand{\diag}{\operatorname{diag}}
\newcommand{\Disc}{\operatorname{Disc}}
\newcommand{\PV}{\operatorname{PV}}
\newcommand{\Hil}{\mathcal H}
\newcommand{\Phys}{\mathcal H_{\mathrm{phys}}}
\newcommand{\phys}{\mathrm{phys}}
\newcommand{\Alg}{\mathcal A}
\newcommand{\Rep}{\mathcal V}
\newcommand{\Path}{\mathsf{Path}}
\newcommand{\LF}{\mathrm{LF}}
\newcommand{\TMD}{\mathrm{TMD}}
\newcommand{\QCD}{\mathrm{QCD}}
\newcommand{\GTMD}{\mathrm{GTMD}}
\newcommand{\GPD}{\mathrm{GPD}}
\newcommand{\Amp}{\mathbf{Amp}}
\newcommand{\Dens}{\mathbf{Dens}}
\newcommand{\Match}{\mathbf{Match}}
\newcommand{\Red}{\mathbf{Red}}
\newcommand{\Proc}{\mathbf{Proc}}
\newcommand{\cO}{\mathcal O}
\newcommand{\cM}{\mathcal M}
\newcommand{\cR}{\mathcal R}
\newcommand{\cS}{\mathcal S}
\newcommand{\cV}{\mathcal V}
\newcommand{\cW}{\mathcal W}
\newcommand{\cE}{\mathcal E}
\newcommand{\cG}{\mathcal G}
\newcommand{\cK}{\mathcal K}
\newcommand{\cQ}{\mathcal Q}
\newcommand{\cP}{\mathcal P}
\newcommand{\cC}{\mathcal C}
\newcommand{\cZ}{\mathcal Z}
\newcommand{\cD}{\mathcal D}
\newcommand{\cN}{\mathcal N}
\newcommand{\cU}{\mathcal U}
\newcommand{\cF}{\mathcal F}
\newcommand{\cA}{\mathcal A}
\newcommand{\cB}{\mathcal B}
\newcommand{\cL}{\mathcal L}
\newcommand{\cJ}{\mathcal J}
\newcommand{\cI}{\mathcal I}
\newcommand{\eps}{\varepsilon}
\newcommand{\etaW}{\eta_{\!W}}
\newcommand{\bbone}{\bm 1}
\newcommand{\abs}[1]{\left|#1\right|}
\newcommand{\norm}[1]{\left\lVert#1\right\rVert}
\newcommand{\inner}[2]{\left\langle #1\middle|#2\right\rangle}
\newcommand{\avg}[1]{\left\langle #1\right\rangle}
\newcommand{\comm}[2]{\left[#1,#2\right]}
\newcommand{\acom}[2]{\left\{#1,#2\right\}}
\newcommand{\Tord}{\mathcal P}
\newcommand{\crossT}[2]{\bigl(#1\!\times\!#2\bigr)_z}
\newcommand{\straight}{\mathrm{str}}
\newcommand{\staple}{\mathrm{stp}}
\newcommand{\unsub}{\mathrm{unsub}}
\newcommand{\subtr}{\mathrm{sub}}
\newcommand{\ren}{\mathrm{ren}}
\newcommand{\odd}{\mathrm{odd}}
\newcommand{\even}{\mathrm{even}}
\newcommand{\soft}{\mathrm{soft}}
\newcommand{\micro}{\mathrm{micro}}
\newcommand{\hard}{\mathrm{hard}}
\newcommand{\rap}{\mathrm{rap}}
\newcommand{\UV}{\mathrm{UV}}
\newcommand{\IR}{\mathrm{IR}}
\newcommand{\LFTO}{\mathrm{LFTO}}

\newcolumntype{Y}{>{\raggedright\arraybackslash}X}
\newcolumntype{L}[1]{>{\raggedright\arraybackslash}p{#1}}

\title{\vspace{-1.0cm}
{\color{navy}\bfseries Volume III: Dynamical Wilson Lines, Rescattering Phases,\\[2mm]
and Naive-\(T\)-Odd GTMDs}\\[4mm]
\large Gauge transport, eikonal light-front dynamics, absorptive cuts,\\
process reversal, soft subtraction, and quark/gluon color structure}
\author{DeuteronWigner project}
\date{29 July 2026}

\begin{document}
\maketitle

\begin{center}
\begin{minipage}{0.94\textwidth}
\colorbox{softblue}{\parbox{0.96\textwidth}{
\textbf{Status and purpose.}
Volume~I defined the regulated gauge-constrained light-front state space and
microscopic mass operator.  Volume~II constructed the common zero-rescattering
quark, antiquark, and gluon GTMD overlaps from those states.  The present
volume supplies the missing process-dependent gauge transport and absorptive
dynamics.  It separates the geometric Wilson path, its eikonal pole
prescription, light-front intermediate-state cuts, rapidity regulation,
soft-overlap subtraction, and process-factorization status.  Its primary
pilot is the leading-twist, zero-skewness nucleon correlator, but every object
is target-spin agnostic so that the same engine can be consumed by the
spin-1 nuclear construction in Volume~IV.  The central output is an
operator-valued rescattering kernel with explicit flavor, link, color,
truncation, and phase provenance; a universal scalar phase multiplying
unrelated TMDs is excluded by construction.
}}
\end{minipage}
\end{center}

\begin{abstract}
Naive-time-reversal-odd transverse-momentum-dependent distributions cannot be
generated by a real zeroth-order light-front overlap alone.  Their leading
power contribution is encoded in process-dependent Wilson lines and the
absorptive support of initial- or final-state rescattering
\cite{BrodskyHwangSchmidt2002,Collins2002,BelitskyJiYuan2003}.  This volume
constructs that dynamics as a typed extension of the common GTMD overlap
formalism.

We fix the covariant-derivative and Fourier conventions, define bare Wilson
transport as a functor from an admissible path groupoid to color fibers, and
construct future- and past-pointing staples with explicit transverse closure.
For a semi-infinite segment the path integral produces the eikonal denominator
\(1/(v\!\cdot\!\ell-i0\eta)\), whose distributional imaginary part changes
sign under \(\eta\to-\eta\).  The pole is not by itself the complete physical
phase: it is embedded in a light-front time-ordered amplitude with
intermediate-state energy denominators, constrained-field and transverse-link
terms, regulator data, and a cut-provenance ledger that forbids counting the
same on-shell support twice.  The one-gluon correction is therefore an
operator kernel in Fock, color, helicity, and orbital space, not an adjustable
complex number.

The formalism distinguishes Hermitian path inversion from antiunitary
future/past reversal, defines link-even and link-odd GTMD components at
nonzero transfer, and derives the familiar forward sign change only after the
complete time-reversal adapter has acted on momenta and spin.  Quark Sivers
and Boer--Mulders functions are different projections of a shared
rescattering interaction: the former uses target-helicity interference in a
vector operator, whereas the latter uses active-quark helicity interference
in a chiral-odd operator.  Gluon correlators retain an ordered pair of
adjoint links and two orthogonal three-adjoint color tensors,
\(if^{abc}\) and \(d^{abc}\); link topology and color contraction remain
separate identity fields.  Transverse moments are matched to renormalized
quark--gluon and tri-gluon pole operators rather than interpreted as
unregulated integrals.

A properly defined TMD or GTMD additionally requires rapidity regulation,
soft subtraction, ultraviolet renormalization, and matching
\cite{EchevarriaGTMD2016,EchevarriaSoft2016}.  We formulate two admissible
microscopic routes: a boundary-only calculation in which rapidity logarithms
and the Collins--Soper kernel are left to Volume~V, and a joint soft-sector
calculation in which both the boundary and kernel are generated and matched
from the same microscopic theory.  Mixing the two routes is prohibited.
Higher eikonal orders are organized by Dyson and Magnus expansions with
independent Wilson-order convergence.  A process registry carries the
factorization theorem, link map, hard color weights, and Glauber status, and
fails closed in channels where separate-hadron TMD factorization is not
established or is known to fail \cite{CollinsQiu2007,RogersMulders2010}.
The volume concludes with formal propositions, analytic benchmarks, software
interfaces, property-based tests, and acceptance criteria for exporting a
dynamical link-resolved nucleon parent to Volumes~IV and~V.
\end{abstract}

\tableofcontents

\section{Role, scope, and inherited contracts}
\label{sec:role}

\subsection{The missing dynamical layer}

Volume~II exported a zero-rescattering parent
\begin{equation}
 W_{a/N}^{(0)}
 =\langle N'|\cO_a\,U_{\gamma}^{(0)}|N\rangle,
 \qquad U_{\gamma}^{(0)}=\bbone,
\label{eq:volume2-zero-parent}
\end{equation}
with complete momentum-fiber, species, flavor, helicity, transverse-rank,
path, color, regulator, scheme, and microscopic-member identity.  For real
state amplitudes and real operator kernels, every path-odd projection of
\cref{eq:volume2-zero-parent} vanishes.  This is a controlled boundary, not a
claim that physical Sivers, Boer--Mulders, or gluon naive-\(T\)-odd functions
are zero.

The present volume replaces the identity transport by
\begin{equation}
 U_\gamma^{(R)}
 =\Tord\exp\!\left[-ig\int_\gamma \dd\xi^\mu
 A_\mu^a(\xi)T_R^a\right],
\label{eq:wilson-transport-intro}
\end{equation}
using the same gauge field, color coupling, state sectors, and
renormalization datum that enter the microscopic Hamiltonian.  The
calculation therefore follows
\begin{equation}
 \left\{\cM_{\LF}^2,|\Psi_N\rangle,\cO_a\right\}
 \longrightarrow
 \left\{U_\gamma,\cK_{\mathrm{resc}},\Disc\right\}
 \longrightarrow
 W_{a/N}^{[\gamma]}
 \longrightarrow
 \left\{F_{\even},F_{\odd}\right\}.
\label{eq:volume3-chain}
\end{equation}

\begin{principle}[Shared interaction, distinct projections]
\label{prin:shared-rescattering}
A quark--gluon vertex, eikonal exchange, intermediate state, or color tensor
is represented once and contributes to every operator projection that uses
it.  Sivers, Boer--Mulders, and gluon naive-\(T\)-odd functions may be
correlated because they share microscopic interactions, but none may be
defined as a scalar multiple of another unless a derived approximation with
an explicit domain and error is supplied.
\end{principle}

\subsection{Primary domain}

The primary pilot retains:
\begin{itemize}
\item a spin-\(\tfrac12\) nucleon target and the accepted Volume~II momentum
      fibers;
\item zero skewness, positive \(x\), and the DGLAP region;
\item leading-twist quark, antiquark, and gluon operators;
\item the first nonvanishing eikonal/Wilson order, followed by a declared
      sequence in Wilson order;
\item a regulator and soft-subtraction contract that can be matched to the
      common TMD scheme of Volume~V.
\end{itemize}
The path and color engine is target-spin independent.  Volume~IV may therefore
apply the same operators to a spin-1 nuclear state without introducing a new
phenomenological phase layer.

\subsection{Explicit nonclaims}

This volume does not claim:
\begin{itemize}
\item that one-gluon exchange is the complete Wilson-line dynamics;
\item that a finite Wilson expansion is gauge, regulator, or scheme
      independent before matching;
\item that every process admits a factorization formula with separately
      defined TMDs;
\item that link topology alone determines a numerical phase;
\item that every imaginary part of a raw helicity amplitude is a naive-
      \(T\)-odd distribution;
\item that Sivers or Boer--Mulders functions are topological invariants;
\item that a phenomenological lensing relation between a TMD and a GPD is an
      operator identity;
\item that a finite-order subtracted correlator must satisfy every exact
      pointwise positivity inequality.
\end{itemize}

\subsection{Interfaces with the volume series}

\begin{longtable}{L{0.17\textwidth}L{0.31\textwidth}L{0.43\textwidth}}
\toprule
Interface & Received or exported object & Contract\\
\midrule
\endhead
Volume I & Regulated states, gauge constraints, Hamiltonian vertices,
current operators, truncation labels & Rescattering uses the same gauge field,
couplings, counterterms, zero-mode prescription, and microscopic-member
identity.\\
Volume II & Typed quark, antiquark, and gluon overlap operators at nonzero
\(\DT\) with \(U^{(0)}=1\) & Recoil, normalization, projectors, and
operator fibers are immutable.  Volume III adds path dynamics but does not
rebuild the parent.\\
Volume IV & Link-resolved nucleon helicity matrices & Nuclear composition
acts on the complete parent and may add nuclear rescattering only through a
separate, matched many-body operator.\\
Volume V & Bare/unsubtracted and microscopic-subtracted boundary objects,
rapidity/soft metadata, Wilson-order manifest & Volume V performs the final
regulator-to-QCD matching, common evolution, and process-level \(W+Y\)
assembly.\\
C1/C2 software spine & Coordinate, rank, sector, path, operator, map-class,
reduction, and provenance identities & No new untyped array route or implicit
phase alternative may bypass these identities.\\
\bottomrule
\end{longtable}

\section{Gauge convention and bare path geometry}
\label{sec:gauge-path}

\subsection{Gauge-field convention}

We fix
\begin{equation}
 D_\mu=\partial_\mu+igA_\mu,
 \qquad
 A_\mu=A_\mu^aT_R^a,
\label{eq:covariant-derivative}
\end{equation}
so that
\begin{equation}
 [D_\mu,D_\nu]=igF_{\mu\nu},
 \qquad
 F_{\mu\nu}
 =\partial_\mu A_\nu-\partial_\nu A_\mu
 +ig[A_\mu,A_\nu].
\label{eq:field-strength-convention}
\end{equation}
For a gauge transformation \(V_R(x)\),
\begin{align}
 \psi_R(x)&\longmapsto V_R(x)\psi_R(x),
\nonumber\\
 A_\mu(x)&\longmapsto
 V_R(x)A_\mu(x)V_R^\dagger(x)
 +\frac{i}{g}\bigl(\partial_\mu V_R(x)\bigr)V_R^\dagger(x).
\label{eq:gauge-transformation}
\end{align}
The fundamental generators obey
\begin{equation}
 \tr(T_F^aT_F^b)=T_F\delta^{ab},
 \qquad T_F=\frac12,
\label{eq:fundamental-normalization}
\end{equation}
and the adjoint generators are
\begin{equation}
 (T_A^a)_{bc}=-if^{abc}.
\label{eq:adjoint-generator}
\end{equation}
Every implementation stores the generator normalization; it is not inferred
from a species label.

\subsection{Oriented paths and endpoint fibers}

An oriented piecewise smooth path is
\begin{equation}
 \gamma:[0,1]\longrightarrow M,
 \qquad
 \gamma(0)=a,
 \quad
 \gamma(1)=b.
\label{eq:path-definition}
\end{equation}
The transport in representation \(R\) is
\begin{equation}
 U_\gamma^{(R)}[b,a]
 =\Tord\exp\!\left[-ig\int_0^1\dd s\,
 \dot\gamma^\mu(s)A_\mu^a(\gamma(s))T_R^a\right].
\label{eq:path-transport}
\end{equation}
It transforms as
\begin{equation}
 U_\gamma^{(R)}[b,a]
 \longmapsto
 V_R(b)U_\gamma^{(R)}[b,a]V_R^\dagger(a).
\label{eq:path-gauge-covariance}
\end{equation}
Thus a Wilson line is a map between endpoint color fibers, not an ordinary
matrix with no source or target.

\begin{definition}[Bare path identity]
\label{def:bare-path-id}
The bare path identity is
\begin{equation}
 \mathfrak g_{\rm bare}
 =\bigl(
 a,b,\gamma,\mathrm{segments},\mathrm{orientation},R,
 \mathrm{endpoint\ closure},\mathrm{Fourier\ sign}
 \bigr).
\label{eq:bare-path-identity}
\end{equation}
It contains geometric and representation data only.  Rapidity regulator,
soft prescription, scales, and renormalization scheme belong to the decorated
operator identity in \cref{sec:operator-hierarchy}.
\end{definition}

\subsection{Composition, inversion, and unitarity}

For composable paths \(\gamma_1:a\to b\) and \(\gamma_2:b\to c\),
\begin{equation}
 U_{\gamma_2\circ\gamma_1}[c,a]
 =U_{\gamma_2}[c,b]U_{\gamma_1}[b,a].
\label{eq:path-composition}
\end{equation}
For Hermitian gauge fields and a real contour,
\begin{equation}
 U_{\gamma^{-1}}[a,b]
 =U_\gamma^\dagger[b,a],
 \qquad
 U_\gamma^\dagger U_\gamma=\bbone.
\label{eq:path-inversion}
\end{equation}

\begin{proposition}[Bare transport functor]
\label{prop:path-functor}
At fixed regulator and finite-dimensional color representation, the
assignment
\begin{equation}
 \cU_R:\Path(M)\longrightarrow\mathrm{ColorFib}_R,
 \qquad
 \gamma\longmapsto U_\gamma^{(R)}
\label{eq:path-functor}
\end{equation}
respects identities, path concatenation, and inversion.
\end{proposition}

\begin{proof}
The path-ordered exponential is the unique solution of the parallel-transport
initial-value problem.  Evolution over adjacent intervals composes, the
zero-length path has unit transport, and reversing the parameter changes the
sign and order of the integral, giving \cref{eq:path-inversion}.
\end{proof}

\subsection{Staple paths and transverse closure}

For a bilocal operator with endpoints \(a=-z/2\) and \(b=z/2\), define a
staple
\begin{equation}
 \gamma_{\eta,v,c_T}
 =\gamma_{a\to a+\eta\infty v}
 \circ\gamma_{\perp,\infty}
 \circ\gamma_{b+\eta\infty v\to b},
 \qquad \eta\in\{+1,-1\},
\label{eq:staple-path}
\end{equation}
where \(v\) is the rapidity direction and \(c_T\) specifies the transverse
closure.  Future and past staples are distinct paths.  Their endpoints are
the same, but they are not identified by path homotopy in a gauge field with
nonzero curvature.  The transverse connector at light-cone infinity is part
of the operator and is essential in gauges where the transverse gauge field
does not vanish asymptotically \cite{BelitskyJiYuan2003}.

The path label stores separately:
\begin{equation}
 \mathrm{direction}=\eta,
 \qquad
 \mathrm{rapidity\ vector}=v,
 \qquad
 \mathrm{closure}=c_T,
 \qquad
 \mathrm{boundary\ condition}=B_\infty.
\label{eq:staple-fields}
\end{equation}
No default closure is inferred from \(\eta\).

\subsection{Path deformation and curvature}

For fixed endpoints, an infinitesimal deformation \(\delta x^\nu(s)\) gives
schematically
\begin{equation}
 \delta U_\gamma[b,a]
 =-ig\int_0^1\dd s\,
 U_{\gamma_{>s}}[b,x_s]
 F_{\mu\nu}(x_s)
 U_{\gamma_{<s}}[x_s,a]
 \dot x^\mu(s)\delta x^\nu(s),
\label{eq:path-deformation}
\end{equation}
with additional endpoint terms when the endpoints move.  Equation
\eqref{eq:path-deformation} makes the role of geometry precise: curvature,
not homotopy alone, controls the physical path dependence.

\begin{requirement}[No topological substitution]
\label{req:no-topological-substitution}
Two paths may be identified only by an exact operator identity, a proven
factorization equivalence, or a controlled approximation with an error
estimate.  Equal connectivity or homotopy class is insufficient in generic
QCD where \(F_{\mu\nu}\neq0\).
\end{requirement}

\section{Operator hierarchy and complete identity}
\label{sec:operator-hierarchy}

\subsection{Four distinct objects}

The calculation distinguishes
\begin{equation}
 \cO_{\rm bare}^{[\gamma]},
 \qquad
 W_{\unsub}^{[\gamma,\cR]},
 \qquad
 W_{\subtr}^{[\gamma,\cR,\cS]},
 \qquad
 W_{\ren}^{[\gamma]}(\mu,\zeta;\mathfrak s).
\label{eq:four-operator-levels}
\end{equation}
The bare bilocal contains fields and geometric transport.  The unsubtracted
matrix element is evaluated in a regulated microscopic state.  The
soft-subtracted object removes the overlap with the declared soft sector.
The renormalized object additionally carries ultraviolet and rapidity
renormalization and belongs to a named TMD/GTMD scheme.

The typed route is
\begin{equation}
\begin{tikzcd}[column sep=large]
W_{\LF}^{(0)}
 \arrow[r,"\mathsf{Wilson}_{\gamma}^{(N_W)}"]
& W_{\unsub}^{[\gamma]}
 \arrow[r,"\mathsf{Sub}_{\soft}"]
& W_{\subtr}^{[\gamma]}
 \arrow[r,"\cZ_{\LF\to\QCD}"]
& W_{\ren}^{[\gamma]}(\mu,\zeta)
 \arrow[r,"\Red"]
& \{F_A\}.
\end{tikzcd}
\label{eq:typed-wilson-route}
\end{equation}
Each arrow has a different mathematical and physical status.

\subsection{Decorated operator identity}

\begin{definition}[Volume III operator identity]
\label{def:volume3-operator-id}
A production operator has identity
\begin{equation}
 \mathfrak o_{\rm III}
 =\Bigl(
 a,J,P_{\rm in},P_{\rm out},\gamma,R,
 [U,U']_g,c_g,
 \cR_{\UV},\cR_{\rap},\cR_{\IR},
 \cS_{\soft},\mu,\zeta,\mathfrak s,
 \mathfrak r,\mathfrak f,\mathfrak m,
 N_W,\mathfrak p
 \Bigr),
\label{eq:volume3-operator-id}
\end{equation}
where \(a\) is species/flavor, \(J\) is the Dirac or Lorentz projection,
\(\gamma\) is the complete path, \([U,U']_g\) is the ordered gluon link pair,
\(c_g\in\{f,d,\mathrm{NA}\}\) is color-contraction status,
\(\mathfrak r\) is transverse-rank identity, \(\mathfrak f\) is the incoming
and outgoing momentum-fiber identity, \(\mathfrak m\) is microscopic member,
\(N_W\) is Wilson/eikonal truncation, and \(\mathfrak p\) is phase provenance.
\end{definition}

The C1/C2 identity is therefore extended, not replaced.  Missing path, pole,
soft, or color data may be recorded as incomplete during migration, but an
operation requiring those fields must fail closed.

\subsection{Hermitian inversion is not process reversal}

Hermitian conjugation acts on an operator by reversing endpoint order and the
geometric path:
\begin{equation}
 \bigl(\cO_{\Gamma}^{[\gamma]}[b,a]\bigr)^\dagger
 =\chi_H(\Gamma)\,
 \cO_{\Gamma_H}^{[\gamma^{-1}]}[a,b].
\label{eq:operator-hermitian-conjugation}
\end{equation}
After the endpoint relabeling and Fourier-variable transformation used in the
GTMD definition, this yields the Volume~II Hermiticity relation at fixed
staple orientation.  Antiunitary time reversal instead maps future to past
staples and transforms spin and momenta.  The two operations must never share
one Boolean field.

\begin{requirement}[Separate involutions]
\label{req:separate-involutions}
The implementation provides separate maps
\begin{equation}
 \mathsf H:\gamma\mapsto\gamma^{-1},
 \qquad
 \mathsf T:\gamma_\eta\mapsto\gamma_{-\eta},
\label{eq:separate-involution-maps}
\end{equation}
with separate tests and convention adapters.
\end{requirement}

\subsection{Process-to-link map}

A process record does not merely name a final state.  It is
\begin{equation}
\begin{aligned}
 \mathfrak P=\bigl(&
 \mathrm{process},\mathrm{measurement},\mathrm{power},
 \mathrm{partonic\ channel},\gamma,[U,U']_g,\\
 &\{C_c^{\hard}\},\mathrm{factorization\ status},
 \mathrm{Glauber\ status},\mathrm{reference}\bigr).
\end{aligned}
\label{eq:process-record}
\end{equation}
The map
\begin{equation}
 \mathsf L_{\Proc}:\mathfrak P\longmapsto
 \bigl(\gamma,[U,U']_g,\{C_c^{\hard}\}\bigr)
\label{eq:process-link-map}
\end{equation}
is defined only when the corresponding factorization statement is adequate
for the requested use.  A process with unknown or broken TMD factorization
cannot silently borrow the SIDIS or Drell--Yan staple.

\section{Wilson expansion and truncation theory}
\label{sec:wilson-expansion}

\subsection{Dyson/path-ordered series}

Parameterize a path by \(x(s)\), \(s\in[0,1]\), and define
\begin{equation}
 B(s)=-ig\dot x^\mu(s)A_\mu^a(x(s))T_R^a.
\label{eq:path-generator}
\end{equation}
Then
\begin{align}
 U_\gamma
 &=\bbone+\sum_{n=1}^{\infty}U_\gamma^{(n)},
\label{eq:dyson-series}\\
 U_\gamma^{(n)}
 &=\int_{0<s_n<\cdots<s_1<1}
 \dd s_1\cdots\dd s_n\,
 B(s_1)\cdots B(s_n).
\label{eq:dyson-term}
\end{align}
The ordered color product is physical.  Replacing it by a scalar
\(C_R^n\) before all contractions are known erases commutator and
\(f/d\)-sensitive information.

\subsection{Fock-space action}

When the gauge field is expanded in creation and annihilation operators,
\(U_\gamma^{(n)}\) contains blocks that change explicit gluon number by up to
\(n\), together with contractions internal to the operator.  In the Volume~I
Fock basis,
\begin{equation}
 P_\nu U_\gamma^{(n)}P_\mu\neq0
\label{eq:wilson-fock-block}
\end{equation}
only when the field content and quantum numbers permit the corresponding
transition.  Thus a Wilson expansion may connect the Volume~II diagonal
zeroth-order overlap to explicit \(qqqg\), \(qqqgg\), or induced-operator
blocks.  Every nonzero block records whether the gluons are explicit state
constituents, operator insertions, or integrated-out modes.

\begin{requirement}[No explicit/induced duplication]
\label{req:no-explicit-induced-duplication}
If a Wilson-induced higher-sector contribution is represented explicitly in
the Fock space, the corresponding effective operator generated by integrating
out that sector is excluded, unless a subtraction two-cell states the overlap
and removes it.
\end{requirement}

\subsection{Magnus representation}

The same transport may be written
\begin{equation}
 U_\gamma=\exp\Omega_\gamma,
 \qquad
 \Omega_\gamma=\Omega_1+\Omega_2+\cdots,
\label{eq:magnus-expansion}
\end{equation}
with
\begin{align}
 \Omega_1&=\int_0^1\dd s_1\,B(s_1),
\label{eq:magnus-1}\\
 \Omega_2&=\frac12\int_0^1\dd s_1\int_0^{s_1}\dd s_2\,
 [B(s_1),B(s_2)].
\label{eq:magnus-2}
\end{align}
For Hermitian gauge fields, a consistently truncated \(\Omega\) is
anti-Hermitian, so \(\exp\Omega^{(\le N)}\) is unitary.  A Dyson truncation is
strictly order-by-order but need not be exactly unitary at finite order.  Both
representations are useful and must agree through their common perturbative
order \cite{Magnus1954}.

\subsection{Independent truncation axes}

The Wilson calculation carries at least
\begin{equation}
 \nu_W=
 \bigl(N_W,N_{\rm web},N_{\rm Fock},N_{\rm basis},
 \Lambda_{\UV},\Lambda_{\IR},\rho_{\rap},N_{\rm quad}\bigr).
\label{eq:wilson-truncation-tuple}
\end{equation}
Here \(N_W\) is field-insertion order, \(N_{\rm web}\) is connected-web or
Magnus order when used, and the remaining labels are inherited state and
numerical truncations.  Convergence in \(N_W\) does not compensate for an
unconverged Fock basis, and regulator variation is not a substitute for
quadrature refinement.

\begin{definition}[Wilson-order residual]
\label{def:wilson-residual}
For a matched observable or operator norm,
\begin{equation}
 \eps_W^{(N)}
 =\frac{\norm{\cR_{N\to N-1}[W^{(\le N)}]-W^{(\le N-1)}}}
 {\max\!\left(\norm{W^{(\le N-1)}},W_{\rm floor}\right)}.
\label{eq:wilson-order-residual}
\end{equation}
The comparison map \(\cR_{N\to N-1}\) removes basis or scheme differences
before the residual is formed.
\end{definition}

\section{Semi-infinite eikonal lines and pole prescriptions}
\label{sec:eikonal-poles}

\subsection{Derivation from the path integral}

Use the field Fourier convention
\begin{equation}
 A_\mu(x)=\int\frac{\dd^d\ell}{(2\pi)^d}
 e^{-i\ell\cdot x}A_\mu(\ell).
\label{eq:gauge-field-fourier}
\end{equation}
For the segment \(x+\eta\lambda v\), \(\lambda\in[0,\infty)\),
\begin{equation}
 \int_0^\infty\dd\lambda\,\eta v^\mu
 e^{-i\eta\lambda v\cdot\ell-\epsilon\lambda}
 =-\frac{i v^\mu}{v\cdot\ell-i0\eta}.
\label{eq:semi-infinite-integral}
\end{equation}
The orientation dependence is therefore encoded in
\begin{equation}
 D_\eta(\ell;v)
 =\frac{1}{v\cdot\ell-i0\eta}.
\label{eq:eikonal-denominator}
\end{equation}
With the exponent convention in \cref{eq:path-transport}, the complete
one-gluon vertex includes the corresponding factor \(-gT_R^av^\mu
D_\eta\), up to the declared momentum-flow convention.

\begin{remark}[Convention identity]
Changing the sign in the Fourier transform, reversing all momentum flow, or
using \(D_\mu=\partial_\mu-igA_\mu\) changes intermediate signs.  The physical
future/past relation is unaffected when every convention is transformed
consistently.  The software therefore stores the Fourier sign, path
orientation, exponent sign, and momentum-flow direction in the pole identity.
\end{remark}

\subsection{Principal value and absorptive support}

The distribution identity is
\begin{equation}
 \frac{1}{x-i0\eta}
 =\PV\frac1x+i\eta\pi\delta(x).
\label{eq:eikonal-distribution-identity}
\end{equation}
Hence
\begin{equation}
 \Im D_\eta(\ell;v)
 =\eta\pi\delta(v\cdot\ell),
 \qquad
 \Im D_{-\eta}=-\Im D_\eta.
\label{eq:eikonal-imaginary}
\end{equation}
The delta function identifies the gluonic-pole surface.  It does not by
itself prove that a named TMD is nonzero: the color, spin, OAM, spectator,
and remaining denominator structure may make the projected interference
vanish.

\begin{proposition}[Orientation-odd pole]
\label{prop:orientation-odd-pole}
For fixed Fourier and momentum-flow conventions, the absorptive part of a
single semi-infinite eikonal denominator is odd under path reversal,
\begin{equation}
 \operatorname{Abs}D_{-\eta}=-\operatorname{Abs}D_\eta.
\label{eq:orientation-odd-pole-result}
\end{equation}
\end{proposition}

\begin{proof}
Apply \cref{eq:eikonal-distribution-identity}.  The principal-value part is
independent of \(\eta\), whereas the delta-function coefficient is linear in
\(\eta\).
\end{proof}

\subsection{Rapidity regulation}

For a lightlike \(v\), loop integrations involving \cref{eq:eikonal-denominator}
contain rapidity divergences not regulated by ordinary dimensional
regularization.  The operator identity must therefore include one of:
\begin{itemize}
\item a non-lightlike direction \(v^2\neq0\);
\item a \(\delta\)-type shift of eikonal denominators;
\item an analytic or exponential rapidity regulator;
\item another regulator with a proven matching to the selected TMD scheme.
\end{itemize}
We denote the regulated denominator by
\begin{equation}
 D_{\eta}^{[\cR_{\rap}]}(\ell;v,\rho_{\rap}).
\label{eq:regulated-eikonal-denominator}
\end{equation}
The rapidity parameter \(\rho_{\rap}\) is not an infrared mass and may not be
used as one.

\begin{requirement}[Rapidity-regulator closure]
\label{req:rapidity-regulator-closure}
Every numerical phase result must be accompanied by a test that the
\(\rho_{\rap}\) dependence cancels or evolves according to the declared
rapidity-renormalization equation after soft subtraction and matching.
\end{requirement}

\subsection{Gauge completion at light-cone infinity}

In light-cone gauge, setting the longitudinal segment to unity does not in
general remove the Wilson-line physics.  Residual gauge fields and the
transverse connector at infinity contribute at leading power
\cite{BelitskyJiYuan2003}.  In a Hamiltonian implementation, the gauge-complete
one-gluon set may include:
\begin{enumerate}
\item exchange from the longitudinal eikonal segment;
\item exchange from the transverse closure;
\item instantaneous or constrained-field terms required by the chosen
      light-front gauge;
\item boundary/zero-mode contributions fixed by the Volume~I prescription.
\end{enumerate}
A gauge-dependent subset is a benchmark, not a production Wilson kernel.

\section{Light-front time-ordered rescattering amplitudes}
\label{sec:lf-rescattering}

\subsection{Resolvent form}

Let \(P_0^-\) be the regulated free or reference light-front Hamiltonian and
\begin{equation}
 G_0(E)=\frac{1}{E-P_0^-+i0}.
\label{eq:lf-resolvent}
\end{equation}
At first nontrivial rescattering order, an operator amplitude contains
pre- and post-insertion terms,
\begin{equation}
 \cA_{\eta}^{(1)}
 =\langle f|V_\eta G_0(E_i)\cO|i\rangle
 +\langle f|\cO G_0(E_i)V_\eta|i\rangle
 +\cA_{\eta,\rm inst}^{(1)}
 +\cA_{\eta,\infty}^{(1)},
\label{eq:first-rescattering-amplitude}
\end{equation}
where \(V_\eta\) is the eikonal interaction derived from the same gauge
sector as the Hamiltonian.  Process kinematics select which ordering is
leading and how it is reorganized into the operator Wilson line.

In an intermediate-state basis,
\begin{equation}
 \cA_{\eta,\rm post}^{(1)}
 =\sum_X
 \frac{\langle f|V_\eta|X\rangle
       \langle X|\cO|i\rangle}
      {E_i-P_X^-+i0},
\label{eq:intermediate-state-amplitude}
\end{equation}
with an analogous pre-insertion term.  The state \(X\) retains Fock sector,
parton identities, color, helicity, OAM, momentum, and regulator labels.

\subsection{Intermediate-state cut}

The light-front energy denominator satisfies
\begin{equation}
 \frac{1}{E-P_X^-+i0}
 =\PV\frac{1}{E-P_X^-}-i\pi\delta(E-P_X^-).
\label{eq:lf-energy-distribution}
\end{equation}
Thus
\begin{equation}
 \Disc\cA^{(1)}
 =-2\pi i\sum_X
 \delta(E_i-P_X^-)
 \langle f|V_\eta|X\rangle
 \langle X|\cO|i\rangle
 +\Disc_{\rm eik}+\Disc_{\rm other},
\label{eq:rescattering-discontinuity}
\end{equation}
where \(\Disc_{\rm eik}\) denotes eikonal-pole support not already represented
by the same fast propagator cut.

\begin{definition}[Cut provenance]
\label{def:cut-provenance}
Each absorptive contribution has
\begin{equation}
 \mathfrak p_{\rm cut}
 =\bigl(
 \mathrm{source},X,\mathrm{denominator},\mathrm{pole},
 \mathrm{cut\ surface},\mathrm{ordering},\mathrm{subtraction\ id}
 \bigr),
\label{eq:cut-provenance}
\end{equation}
with source in
\begin{equation}
 \{\mathrm{EIKONAL\ POLE},\mathrm{LF\ ENERGY\ CUT},
 \mathrm{ZERO\ MODE},\mathrm{BOUNDARY},\mathrm{MATCHING},
 \mathrm{HARD\ PHASE}\}.
\label{eq:phase-source-enum}
\end{equation}
\end{definition}
A hard-scattering phase belongs to \(\Proc\), not to the hadron TMD.  A cut
represented by both a full propagator and its eikonal approximation must be
assigned to one route with a subtraction relation.

\subsection{General one-gluon kernel}

For active representation \(R\) and a target/spectator current
\(\cJ_{X\alpha}^{a\mu}\), a general regulated one-gluon kernel is
\begin{equation}
 \cK_{\eta;\alpha'\alpha}^{(1)}
 =g^2\sum_{a,b,X}\int[\dd\ell]\,
 \cJ_{\alpha'X}^{a\mu}(\ell)
 D_{\mu\nu}^{ab}(\ell;\cR_{\UV},\cR_{\IR})
 T_R^b v^\nu
 D_{\eta}^{[\cR_{\rap}]}(\ell;v)
 \cS_{X\alpha}(\ell),
\label{eq:one-gluon-kernel}
\end{equation}
where \(\alpha,\alpha'\) denote complete Fock/helicity/OAM labels and
\(\cS_{X\alpha}\) contains the remaining state overlap, momentum delta
functions, and energy denominators.  Equation \eqref{eq:one-gluon-kernel} is
an operator-valued matrix.  The familiar screened transverse convolution is a
possible benchmark reduction, not the definition of the production kernel.

The corresponding GTMD correction is
\begin{equation}
 \delta W_{\eta,A}^{(1)}
 =\sum_{\alpha',\alpha}
 \psi_{\alpha'}^*\,
 \cP_A\cK_{\eta;\alpha'\alpha}^{(1)}\,
 \psi_{\alpha},
\label{eq:gtmd-kernel-correction}
\end{equation}
where \(\cP_A\) is the Volume~II operator/helicity/rank projector.  Different
named functions use different \(\cP_A\) on the same kernel.

\subsection{Gauge-covariance and Ward tests}

In a covariant gauge, the propagator contains a longitudinal term
proportional to \(\ell_\mu\ell_\nu\).  Gauge-parameter independence requires
the complete current and eikonal set to satisfy the relevant Ward identity,
including instantaneous and boundary terms.  A numerical diagnostic is
\begin{equation}
 \eps_{\rm Ward}
 =\frac{\norm{\ell_\mu\cJ^\mu_{\rm complete}}}
 {\max(\norm{\cJ_{\rm complete}},J_{\rm floor})}.
\label{eq:ward-residual}
\end{equation}
Production acceptance requires \(\eps_{\rm Ward}\) to decrease with basis and
Wilson-order refinement and to lie below the declared tolerance.

For a color-singlet state,
\begin{equation}
 \left(\sum_{i\in\nu}T_i^a\right)|\Psi_N\rangle=0.
\label{eq:color-singlet-charge}
\end{equation}
Hence, for active parton \(j\),
\begin{equation}
 \sum_{s\neq j}T_s^a|\Psi_N\rangle
 =-T_j^a|\Psi_N\rangle.
\label{eq:spectator-color-closure}
\end{equation}
This is an exact color-closure test, not permission to replace all spectator
dynamics by one Casimir factor.

\subsection{Phase is a relative amplitude property}

A single overall phase of the complete state is unobservable.  A naive-
\(T\)-odd projection requires an absorptive difference between amplitude
components that are selected by distinct spin/OAM projectors.  Schematically,
\begin{equation}
 F_{\odd,A}^{[\eta]}
 \propto
 \sum_{\alpha'\neq\alpha}
 \Im\!\left[
 \psi_{\alpha'}^*
 \cK_{\eta;\alpha'\alpha}
 \psi_\alpha
 \right]\cP_A^{\alpha'\alpha}.
\label{eq:relative-phase-structure}
\end{equation}
Multiplying every amplitude by the same \(e^{i\chi}\) leaves
\cref{eq:relative-phase-structure} unchanged.

\begin{proposition}[No scalar-phase generation]
\label{prop:no-scalar-phase}
Let \(\cK_\eta=e^{i\chi_\eta}\bbone\) act identically on every internal
component.  Then it cannot generate a new naive-\(T\)-odd distribution from a
zeroth-order correlator in which that distribution vanishes.
\end{proposition}

\begin{proof}
The scalar phase cancels between amplitude and conjugate amplitude in every
bilinear matrix element.  It changes neither relative spin/OAM phases nor the
projected density matrix.
\end{proof}

\section{Link-even and link-odd correlator components}
\label{sec:link-parity}

\subsection{Antiunitary time-reversal adapter}

Time reversal is antiunitary and acts on states, spin, momenta, field
operators, and paths.  Let \(\Theta\) denote the represented antiunitary map
and \(\mathsf t\) the induced transformation of the kinematic and operator
labels.  At operator level,
\begin{equation}
 \Theta\,
 \cO_{J}^{[\gamma_\eta]}(P',P)
 \Theta^{-1}
 =\chi_T(J)\,
 \cO_{J_T}^{[\gamma_{-\eta}]}
 (\mathsf tP',\mathsf tP),
\label{eq:time-reversal-operator}
\end{equation}
where \(\chi_T(J)\) and the helicity phase convention are stored in the
operator registry.  The transformed matrix element is complex conjugated by
antiunitarity.

Define the complete time-reversal adapter on a GTMD helicity matrix by
\begin{equation}
 \bigl(\mathfrak T W^{[-\eta]}\bigr)_{\Lambda'\Lambda}
 =\chi_T(J)
 \sum_{\bar\Lambda',\bar\Lambda}
 \Theta_{\Lambda'\bar\Lambda'}
 \Theta_{\Lambda\bar\Lambda}^{*}
 \left[
 W_{\bar\Lambda'\bar\Lambda}^{[-\eta]}
 (\mathsf t x,\mathsf t\kT,\mathsf t\DT)
 \right]^*.
\label{eq:time-reversal-adapter}
\end{equation}
The explicit form of \(\Theta_{\Lambda\bar\Lambda}\) depends on the
light-front helicity phase convention; it is generated from the represented
spin basis rather than hard-coded as a universal sign.

\subsection{Path parity at nonzero transfer}

\begin{definition}[Link-even and link-odd parent]
\label{def:link-even-odd}
For a pair of factorization-compatible future/past paths, define
\begin{align}
 W_{\even}
 &=\frac12\left(W^{[+]}+\mathfrak T W^{[-]}\right),
\label{eq:link-even-parent}\\
 W_{\odd}
 &=\frac12\left(W^{[+]}-\mathfrak T W^{[-]}\right).
\label{eq:link-odd-parent}
\end{align}
\end{definition}
This definition is valid before the forward limit.  It does not identify the
raw imaginary part of one helicity amplitude with a T-odd GTMD.

In a scalar GTMD basis whose explicit factors of \(i\), \(\epsilon_T^{ij}\),
and helicity phases have been fixed, one may write
\begin{equation}
 X_A^{[\eta]}
 =X_{A,\even}+i\eta X_{A,\odd},
\label{eq:scalar-even-odd-decomposition}
\end{equation}
with real coefficient functions in a domain where the standard Hermiticity
and time-reversal relations apply.  The registry records whether the named
basis coefficient includes the explicit \(i\).  Otherwise, the numerical
``imaginary part'' can be convention dependent.

\subsection{Hermiticity at nonzero transfer}

For each scalar basis coefficient \(X_A\), Hermiticity gives
\begin{equation}
 X_A^*(x,\kT^2,\kT\!\cdot\!\DT,\DT^2;\eta)
 =\sigma_H(A)
 X_A(x,\kT^2,-\kT\!\cdot\!\DT,\DT^2;\eta),
\label{eq:gtmd-hermiticity-sign}
\end{equation}
where \(\sigma_H(A)=\pm1\) is determined by the chosen tensor and spin basis
\cite{Meissner2009}.  Thus an imaginary component can be required by
Hermiticity to be odd in \(\kT\cdot\DT\) while still being link even.  This is
why path parity must be computed with \cref{eq:link-even-parent,eq:link-odd-parent}
 rather than inferred from \(\Im X_A\) alone.

\begin{requirement}[Three independent labels]
\label{req:three-phase-labels}
Every complex GTMD coefficient stores separately:
\begin{equation}
 \bigl(
 \sigma_H(A),
 \sigma_T(A),
 \mathrm{explicit\ basis\ phase}
 \bigr).
\label{eq:three-phase-labels}
\end{equation}
Hermiticity parity, time-reversal/link parity, and basis phase may not be
collapsed into one \texttt{t\_odd} flag.
\end{requirement}

\subsection{Zero-rescattering theorem}

\begin{theorem}[Controlled zero of link-odd projections]
\label{thm:zero-rescattering}
Assume:
\begin{enumerate}
\item the Volume~II state amplitudes and zeroth-order operator kernels can be
      chosen real in the declared basis;
\item there is no independent absorptive hard or matching coefficient in the
      parent operator;
\item \(g\to0\) or the Wilson interaction is disabled so
      \(U_{\gamma_+}=U_{\gamma_-}=\bbone\).
\end{enumerate}
Then
\begin{equation}
 W_{\odd}=0
\label{eq:zero-rescattering-odd-zero}
\end{equation}
and every named naive-\(T\)-odd TMD projected from it vanishes.
\end{theorem}

\begin{proof}
Under the assumptions, the future and past operators reduce to the same real
zeroth-order matrix element.  The time-reversal adapter maps that matrix
element to itself, so the difference in \cref{eq:link-odd-parent} is zero.
\end{proof}

\subsection{Forward sign reversal theorem}

At \(\DT=0\), after spin reflection and convention adapters have been
applied, a real named coefficient satisfies
\begin{equation}
 F_{A,\even}^{[+]}(x,\kT)
 =F_{A,\even}^{[-]}(x,\kT),
 \qquad
 F_{A,\odd}^{[+]}(x,\kT)
 =-F_{A,\odd}^{[-]}(x,\kT).
\label{eq:forward-link-reversal}
\end{equation}

\begin{theorem}[SIDIS--Drell--Yan sign reversal]
\label{thm:sidis-dy-sign}
Suppose a quark TMD enters two processes for which standard leading-power TMD
factorization is valid, and the process map assigns the factorization-related
future and past staples.  Then every link-odd quark distribution changes sign
between those process definitions, whereas every link-even quark distribution
is unchanged \cite{Collins2002,JiMaYuanSpin2004}.
\end{theorem}

\begin{remark}
The theorem is an operator statement conditional on the factorization and
path map.  It is not a rule that any phenomenological asymmetry in any two
processes differs only by one sign; hard factors, fragmentation, phase-space
weights, flavor sums, and factorization breaking remain process dependent.
For gluons, ordered link pairs and multiple color structures generally require
more than one sign.
\end{remark}

\section{Quark naive-\texorpdfstring{\(T\)}{T}-odd GTMD and TMD sectors}
\label{sec:quark-todd}

\subsection{Vector and chiral-odd operator projections}

For quarks and antiquarks, the two canonical leading-twist path-odd TMD
families for a spin-\(\tfrac12\) target are generated from different operator
sectors:
\begin{align}
 f_{1T}^{\perp q[\eta]}:
 &\quad J=\gamma^+,
 \quad \mathrm{target}=T,
 \quad \mathrm{active\ quark}=U,
 \quad r=1,
\label{eq:sivers-identity}\\
 h_1^{\perp q[\eta]}:
 &\quad J=i\sigma^{i+}\gamma_5,
 \quad \mathrm{target}=U,
 \quad \mathrm{active\ quark}=T,
 \quad r=1.
\label{eq:boer-mulders-identity}
\end{align}
Their names do not determine their kernels.  The Volume~II projector registry
supplies the exact tensor normalization and reference mass.

\subsection{Sivers interference}

In a target-helicity basis, a transverse target polarization is built from
off-diagonal nucleon helicity elements.  The quark Sivers projection is
schematically
\begin{equation}
 \frac{\epsilon_T^{ij}k_{Ti}S_{Tj}}{M_N}
 f_{1T}^{\perp q[\eta]}
 \;\propto\;
 \sum_{\lambda_q,X}
 \Im\!\left[
 \cA_{\Lambda'=+,\lambda_q,X}^{[\eta]*}
 \cA_{\Lambda=-,\lambda_q,X}^{[\eta]}
 \right],
\label{eq:sivers-helicity-interference}
\end{equation}
up to the stored helicity and epsilon-tensor conventions.  The active quark
projection is helicity diagonal; the target helicity changes by one unit, so
the contributing Fock amplitudes contain the compensating orbital harmonic.

A useful zero test is
\begin{equation}
 \psi_{\Delta L_z=\pm1}=0
 \quad\Longrightarrow\quad
 f_{1T}^{\perp q}=0
\label{eq:sivers-oam-zero}
\end{equation}
for the minimal leading-twist overlap.  The implementation derives the exact
selection rule from total-\(J^z\) equivariance; it does not assume a fixed
sign formula independent of phase convention.

\subsection{Boer--Mulders interference}

The Boer--Mulders function correlates transverse quark spin with \(\kT\) in an
unpolarized target.  Its helicity representation uses active-quark helicity
off-diagonal terms at target-helicity average,
\begin{equation}
 \frac{\epsilon_T^{ij}k_{Tj}}{M_N}
 h_1^{\perp q[\eta]}
 \;\propto\;
 \frac12\sum_{\Lambda,X}
 \Im\!\left[
 \cA_{\Lambda,+,X}^{[\eta]*}
 \cA_{\Lambda,-,X}^{[\eta]}
 \right]_i.
\label{eq:boer-mulders-helicity-interference}
\end{equation}
The quark helicity change is balanced by an orbital harmonic.  Removing the
required helicity-flip or OAM block must make the projection vanish.

\subsection{Shared kernel does not imply proportionality}

Using \cref{eq:gtmd-kernel-correction},
\begin{align}
 f_{1T}^{\perp q}
 &=\cP_{\rm Siv}\!
 \left[\psi^*\cK_\eta\psi\right],
\label{eq:sivers-projector-kernel}\\
 h_1^{\perp q}
 &=\cP_{\rm BM}\!
 \left[\psi^*\cK_\eta\psi\right].
\label{eq:bm-projector-kernel}
\end{align}
Because \(\cP_{\rm Siv}\neq\cP_{\rm BM}\), no operator identity of the form
\begin{equation}
 h_1^{\perp q}=C_q f_{1T}^{\perp q}
\label{eq:forbidden-bm-sivers-identity}
\end{equation}
follows.  A restricted spectator approximation may yield an approximate
ratio, but its error must include omitted helicity, OAM, sea, and Wilson-order
blocks and it may not define the central microscopic model.

\begin{requirement}[Cross-family prediction test]
\label{req:cross-family-prediction}
If a restricted set of Sivers data is used to calibrate shared rescattering
parameters, at least one Boer--Mulders or other link-odd operator family must
remain withheld.  Refitting a separate phase normalization for the withheld
family does not count as prediction.
\end{requirement}

\subsection{Flavor and antiquark structure}

The kernel carries the active flavor, the color/spin content of every
intermediate state, and the sea-sector identity.  Flavor dependence can arise
from:
\begin{itemize}
\item different \(u\) and \(d\) wave-function amplitudes and spectator
      channels;
\item quark masses and chiral interactions;
\item flavor-dependent sea Fock sectors;
\item interference among different OAM and spin-flavor couplers;
\item charge-symmetry-breaking operators where retained.
\end{itemize}
Antiquark Sivers and Boer--Mulders functions are generated from positive-\(x\)
antiquark active slots or a named induced operator.  They are not copied from
the quark result and are not forced to share its sign.

\subsection{Renormalized transverse moment and gluonic pole}

A naive moment
\begin{equation}
 \int\dd^2\kT\,k_T^\alpha\Phi^{[\eta]}(x,\kT)
\label{eq:naive-transverse-moment}
\end{equation}
can be ultraviolet divergent.  The formal object is a renormalized
transverse-derivative or Bessel-weighted limit.  We define its normalization
so that
\begin{equation}
 \Phi_{\partial,\ren}^{\alpha[\gamma]}(x)
 =\widetilde\Phi_{\partial,\ren}^{\alpha}(x)
 +C_G^{[\gamma]}\,
 \Phi_{G,\ren}^{\alpha}(x,x),
\label{eq:quark-gluonic-pole-decomposition}
\end{equation}
where \(\Phi_G\) is the soft-gluon-pole quark--gluon operator and
\(C_G^{[\gamma]}\) is the link-dependent gluonic-pole factor
\cite{BoerMuldersPijlman2003,BuffingMulders2011}.  For the simplest future and
past staples,
\begin{equation}
 C_G^{[+]}=-C_G^{[-]}
\label{eq:quark-gluonic-pole-sign}
\end{equation}
under the chosen normalization.

The Sivers--Qiu--Sterman adapter is
\begin{equation}
 f_{1T}^{\perp(1)q[\gamma]}(x;\mu,\zeta)
 =\cN_{\rm QS}(\mathfrak s,\epsilon_T,\mathrm{Fourier})
 C_G^{[\gamma]}T_F^q(x,x;\mu)
 +\cO(\alpha_s),
\label{eq:sivers-qiu-sterman-adapter}
\end{equation}
where \(\cN_{\rm QS}\) stores all sign, mass, \(\pi\), and normalization
conventions.  The relation is a matched operator statement in the overlap
region of TMD and collinear twist-3 factorization, not a raw-grid integration
identity \cite{JiQiuVogelsangYuan2006,JiQiuVogelsangYuanSIDIS2006}.

\section{Gluon ordered links and independent color structures}
\label{sec:gluon-todd}

\subsection{Two-link operator identity}

The gauge-invariant gluon correlator contains two ordered adjoint transports,
\begin{equation}
 \Gamma^{ij[U,U']}_{g}
 \sim
 \left\langle P'\left|
 F^{+i}_a(a)
 U_{ab}[a,b]
 F^{+j}_b(b)
 U'_{ba}[b,a]
 \right|P\right\rangle.
\label{eq:gluon-two-link-operator}
\end{equation}
The pair \([U,U']\) is ordered.  Interchanging the links, reversing only one
link, or changing a closure can change the operator.  The C1/C2 path identity
therefore stores
\begin{equation}
 \mathfrak g_g
 =\bigl(\gamma_1,R_A;\gamma_2,R_A;
 \mathrm{order},\mathrm{trace\ closure}\bigr).
\label{eq:gluon-link-pair-identity}
\end{equation}

\subsection{Three-adjoint invariant tensors}

For \(SU(N_c\ge3)\), the independent invariant tensors relevant to one
gluonic-pole insertion are
\begin{equation}
 \cC_f^{abc}=if^{abc},
 \qquad
 \cC_d^{abc}=d^{abc}.
\label{eq:f-d-color-tensors}
\end{equation}
They obey
\begin{align}
 f^{abc}d^{abc}&=0,
\label{eq:fd-orthogonality}\\
 f^{abc}f^{abc}&=N_c(N_c^2-1),
\label{eq:ff-norm}\\
 d^{abc}d^{abc}&=\frac{(N_c^2-4)(N_c^2-1)}{N_c}.
\label{eq:dd-norm}
\end{align}
For \(N_c=3\), the last two norms are \(24\) and \(40/3\).

\begin{definition}[Gluon color projection]
\label{def:gluon-color-projection}
Given a three-adjoint kernel \(K^{abc}\), define normalized projections
\begin{align}
 K_f&=\frac{-i f^{abc}K^{abc}}{f^{def}f^{def}},
\label{eq:f-projection}\\
 K_d&=\frac{d^{abc}K^{abc}}{d^{def}d^{def}}.
\label{eq:d-projection}
\end{align}
\end{definition}
No numerical interpolation, basis rotation, or process weight may merge these
orthogonal channels before the process layer.

\subsection{Gluonic-pole decomposition}

The renormalized first transverse moment has the generalized form
\begin{equation}
 \Gamma_{\partial,\ren}^{\alpha[U,U']}(x)
 =\widetilde\Gamma_{\partial,\ren}^{\alpha}(x)
 +\sum_{c\in\{f,d\}}
 C_{G,c}^{[U,U']}\,
 \Gamma_{G,c,\ren}^{\alpha}(x,x),
\label{eq:gluon-gluonic-pole-decomposition}
\end{equation}
where \(C_{G,c}^{[U,U']}\) are process/link-dependent coefficients and
\(\Gamma_{G,f}\), \(\Gamma_{G,d}\) are independent universal operator
matrix elements within the domain of generalized universality
\cite{BuffingGluon2013}.  At higher transverse rank, multiple gluonic poles
can introduce additional color structures and universal functions; the
first-moment two-channel decomposition must not be extrapolated blindly.

A factorized process contribution is therefore
\begin{equation}
 F_{g,\Proc}^{\odd}
 =C_{\hard,f}^{\Proc}F_g^{(f)}
 +C_{\hard,d}^{\Proc}F_g^{(d)}
\label{eq:process-fd-mixture}
\end{equation}
only after the process map supplies the hard color weights.  A default
\(F_g^{(f)}+F_g^{(d)}\) is forbidden.

\subsection{Link topology and color contraction are different labels}

Common simple link pairs include
\begin{equation}
 [+,+],\quad[-,-],\quad[+,-],\quad[-,+],
\label{eq:common-gluon-link-pairs}
\end{equation}
with possible additional Wilson loops.  Same-direction and opposite-direction
pairs are often called Weizs\"acker--Williams-type and dipole-type in
appropriate contexts, but that terminology does not replace the explicit
path.  Likewise, an \(f\)- or \(d\)-type color projection is not globally
identical to one nominal link pair.  The process-dependent relation is stored
as a map, not encoded in an enum alias \cite{Boer2016}.

\begin{requirement}[Independent gluon identities]
\label{req:independent-gluon-identities}
Every gluon naive-\(T\)-odd object carries both
\begin{equation}
 \bigl([U,U'],c_g\bigr),
 \qquad c_g\in\{f,d\}.
\label{eq:gluon-two-label-requirement}
\end{equation}
An object with one field missing is incomplete and may not enter a process
calculation that requires it.
\end{requirement}

\subsection{Small-\texorpdfstring{\(x\)}{x} specialization}

For certain opposite-direction dipole links, the small-\(x\) limit can reduce
the correlator to a Wilson-loop operator, and several transverse-spin gluon
TMDs can share a spin-dependent odderon origin
\cite{BoerEchevarriaMuldersZhou2016,Boer2016}.  This is a specialized
high-energy limit with its own evolution and power counting.  It is recorded
as an additional matching route,
\begin{equation}
 \Gamma_g^{[+,-]}
 \xrightarrow[x\to0]{\mathsf{HE\ match}}
 \cW_{\rm loop}^{\odd},
\label{eq:small-x-wilson-loop-route}
\end{equation}
not as a universal replacement for moderate-\(x\) \(f/d\) dynamics.

\section{Soft subtraction, rapidity matching, and phase budgeting}
\label{sec:soft-subtraction}

\subsection{Why the unsubtracted phase is not yet a TMD}

The microscopic Wilson-line matrix element contains collinear, soft, and
possibly Glauber regions.  A properly defined TMD or GTMD requires a soft
overlap subtraction and rapidity renormalization
\cite{JiMaYuan2004,Collins2011,EchevarriaGTMD2016}.  Schematically,
\begin{equation}
 W_{\ren}^{[\gamma]}(\mu,\zeta)
 =Z_{\UV}^{[\gamma]}(\mu,\cR_{\UV})
 R_{\rap}^{[\gamma]}(\zeta,\cR_{\rap})
 W_{\unsub}^{[\gamma]}(\cR)
 \left[S^{[\gamma]}(\cR)\right]^{-1/2}.
\label{eq:renormalized-tmd-gtmd}
\end{equation}
The placement of the square root and finite factors is scheme dependent; the
operator identity stores the selected convention.

At first order,
\begin{equation}
 W_{\ren}^{(1)}
 =W_{\unsub}^{(1)}
 -\frac12 S^{(1)}W^{(0)}
 +Z_{\UV}^{(1)}W^{(0)}
 +R_{\rap}^{(1)}W^{(0)}.
\label{eq:first-order-subtraction}
\end{equation}
In common schemes the soft factor is spin independent and does not by itself
create a target-spin-odd tensor, but cancellation of regulator and phase
terms must be verified rather than assumed.

\subsection{Phase budget}

\begin{definition}[Phase budget]
\label{def:phase-budget}
For each projected operator family,
\begin{equation}
 \mathfrak B_A
 =\left\{
 \Im W_{\unsub,A},
 \Im W_{\soft,A},
 \Im Z_A,
 \Im R_{\rap,A},
 \Im H_A,
 \mathrm{subtraction\ relations}
 \right\}.
\label{eq:phase-budget}
\end{equation}
\end{definition}
The hadron TMD receives only the subtracted collinear contribution.  A hard
phase \(\Im H_A\) is retained in the process kernel.  A soft/Glauber region
represented in both \(W_{\unsub}\) and a separate soft function is removed by
the named subtraction map.

The production report includes
\begin{equation}
 \eps_{\rm budget}
 =\norm{
 \Im W_{\ren,A}
 -\mathsf{Compose}(\mathfrak B_A)
 }.
\label{eq:phase-budget-residual}
\end{equation}
A small final number without a closed budget is not accepted.

\subsection{Two admissible microscopic routes}

There are two internally consistent strategies.

\paragraph{Route A: boundary-only rescattering.}
The microscopic calculation produces a finite link-dependent boundary at
\((\mu_0,\zeta_0)\), with the logarithmic soft/rapidity overlap subtracted:
\begin{equation}
 W_{\micro,\subtr}^{[\gamma]}(\mu_0,\zeta_0)
 \xrightarrow{\cZ_{\micro\to\TMD}}
 W_{\TMD}^{[\gamma]}(\mu_0,\zeta_0).
\label{eq:boundary-only-route}
\end{equation}
Volume~V supplies the Collins--Soper kernel and evolves away from the boundary.

\paragraph{Route B: joint microscopic soft sector.}
The same microscopic theory computes both the subtracted boundary and the
rapidity kernel,
\begin{equation}
 \left\{W_{\micro}^{[\gamma]},S_{\micro},\cD_{\micro}\right\}
 \xrightarrow{\cZ}
 \left\{W_{\TMD}^{[\gamma]},S_{\QCD},\cD_{\QCD}\right\}.
\label{eq:joint-soft-route}
\end{equation}
This route is more complete but requires a controlled rapidity-renormalization
calculation in the microscopic regulator.

\begin{requirement}[Exclusive soft routes]
\label{req:exclusive-soft-routes}
Route A and Route B are mutually exclusive composition alternatives.  A
microscopic eikonal ladder absorbed into \(\cD_{\micro}\) may not also be
resummed by an external Collins--Soper kernel without an explicit overlap
subtraction.
\end{requirement}

\subsection{Rapidity evolution consistency}

For the matched object,
\begin{equation}
 \frac{\dd}{\dd\ln\sqrt\zeta}
 \ln W_{a}^{[\gamma]}(x,\bM;\mu,\zeta)
 =-\cD_a(\bM;\mu),
\label{eq:cs-evolution-volume3}
\end{equation}
and
\begin{equation}
 \frac{\dd}{\dd\ln\mu}\cD_a(\bM;\mu)
 =\Gamma_{\rm cusp}^a(\alpha_s).
\label{eq:cs-integrability-volume3}
\end{equation}
The kernel is spin independent at leading power for a given color
representation, whereas the link-odd spin structure resides in the boundary
and matching coefficients.  Volume~III exports all data needed to test that
future and past boundaries evolve with the same kernel and retain exact sign
reversal.

\subsection{Scheme covariance}

Let \(\mathfrak s\) and \(\mathfrak s'\) be two valid TMD schemes.  A finite
scheme map acts as
\begin{equation}
 W^{[\gamma]}_{\mathfrak s'}
 =C_{\mathfrak s'\leftarrow\mathfrak s}^{[R]}
 W^{[\gamma]}_{\mathfrak s},
 \qquad
 H_{\mathfrak s'}
 =H_{\mathfrak s}
 \left(C_{\mathfrak s'\leftarrow\mathfrak s}^{[R]}\right)^{-1},
\label{eq:scheme-covariance}
\end{equation}
so the observable is invariant to the stated perturbative order.  A
link-odd sign difference cannot be repaired by changing schemes; a scheme
change must act consistently on both link orientations and the hard factor.

\section{Higher eikonal orders and non-Abelian organization}
\label{sec:higher-orders}

\subsection{Ordered multi-exchange kernel}

The one-gluon result is a benchmark, not a license to replace the Wilson line
by a fitted phase.  At order \(n\), the bare contribution along an oriented
segment is
\begin{align}
 U_{\gamma}^{(n)}
 & =(-ig)^n
 \int_{0<s_n<\cdots<s_1<1}
 \prod_{r=1}^{n}\dd s_r\,
 \dot\gamma^{\mu_r}(s_r)
 A_{\mu_1}^{a_1}(\gamma(s_1))\cdots
 A_{\mu_n}^{a_n}(\gamma(s_n))
 t_R^{a_1}\cdots t_R^{a_n} .
\label{eq:ordered-multigluon}
\end{align}
The matrix order in \cref{eq:ordered-multigluon} is physical.  Even when two
exchanges have identical scalar kernels, interchanging them changes the color
operator whenever \([t_R^{a_1},t_R^{a_2}]\ne0\).  The corresponding
rescattering correction is therefore a map
\begin{equation}
 \cK_{\eta}^{(n)}:
 \Hil_{\nu,\phys}^{(r)}
 \longrightarrow
 \bigoplus_{\mu\in\cF_r}
 \Hil_{\mu,\phys}^{(r)},
\label{eq:multigluon-kernel-map}
\end{equation}
with explicit source and target sectors, not a scalar multiplier.

Let \(\mathfrak X_n\) denote a complete retained intermediate-state history,
including the sector after every exchange.  In light-front time-ordered form,
\begin{align}
 \cK_{\eta}^{(n)}
 ={}&
 \sum_{\mathfrak X_n}
 \cV_{\eta,n}
 \frac{1}{P^-_i-P^-_{X_{n-1}}+i0\sigma_{n-1}}
 \cV_{\eta,n-1}\cdots
 \frac{1}{P^-_i-P^-_{X_1}+i0\sigma_1}
 \cV_{\eta,1},
\label{eq:lf-multiexchange-resolvent}
\end{align}
where each \(\sigma_j\) is fixed by the time ordering and process, and every
\(\cV_{\eta,j}\) carries its own eikonal pole.  A cut may pass through more
than one denominator.  The cut ledger must therefore identify the cut set,
not merely label an entire diagram ``imaginary.''

\begin{requirement}[Complete history identity]
\label{req:complete-history}
Two multi-exchange terms are identical only when their path order, color word,
intermediate sectors, momentum routing, pole prescriptions, regulator labels,
and cut set agree.  Terms that differ in any of these fields may interfere,
but they may not be deduplicated by comparing only their final numerical
arrays.
\end{requirement}

\subsection{Dyson and Magnus truncations are distinct}

The Dyson partial sum
\begin{equation}
 U_{\gamma,[N]}^{\rm D}
 =\sum_{n=0}^{N}U_{\gamma}^{(n)}
\label{eq:dyson-partial-sum}
\end{equation}
and the Magnus partial exponent
\begin{equation}
 U_{\gamma,[N]}^{\rm M}
 =\exp\!\left[\sum_{n=1}^{N}\Omega_n[\gamma]\right]
\label{eq:magnus-partial-sum}
\end{equation}
are different approximations at finite \(N\).  The first is transparent for
fixed-order cuts; the second preserves group-valued transport more naturally
when the retained \(\Omega_n\) are anti-Hermitian.  They become equivalent
only after expansion to a common perturbative order.

\begin{proposition}[Fixed-order equivalence]
\label{prop:dyson-magnus-equivalence}
If \(\Omega_1,\ldots,\Omega_N\) are constructed from the same regulated
connection as \(U^{(1)},\ldots,U^{(N)}\), then
\begin{equation}
 U_{\gamma,[N]}^{\rm M}
 -U_{\gamma,[N]}^{\rm D}
 =\cO(g^{N+1})
\label{eq:dyson-magnus-difference}
\end{equation}
after both sides are expanded in \(g\).  The statement does not imply that
unexpanded numerical partial sums agree beyond the controlled order.
\end{proposition}

\begin{proof}
The Magnus coefficients are defined recursively so that the exponential
solves the same first-order transport equation as the path-ordered Dyson
series.  Matching equal powers of \(g\) fixes \(\Omega_n\) through order
\(n\); the first unmatched term is therefore order \(g^{N+1}\)
\cite{Magnus1954}.
\end{proof}

The implementation records
\begin{equation}
 \mathfrak t_{\rm Wilson}
 =(	ext{representation},N,\text{Dyson/Magnus},
   \text{expanded/unexpanded},\text{cut order},\text{regulator}),
\label{eq:wilson-truncation-id}
\end{equation}
so an unexpanded exponent cannot be compared with a strict fixed-order result
without an explicit expansion adapter.

\subsection{Non-Abelian exponentiation and connected color structure}

For vacuum soft factors and related Wilson-line correlators, the logarithm is
organized by connected color structures or webs rather than by an arbitrary
exponential of the one-gluon graph \cite{Gatheral1983,FrenkelTaylor1984}.
This fact is useful but must not be overextended.  The open, hadron-matrix-
element Wilson operator in the present calculation contains state-dependent
insertions and may mix Fock sectors.  We therefore write schematically
\begin{equation}
 \ln \avg{U_{\gamma_1}\cdots U_{\gamma_m}}
 =\sum_{w\in\mathfrak W}
 C_w\,F_w,
\label{eq:web-expansion}
\end{equation}
only for an object for which the web theorem and its color-connected
hypotheses have been established.

\begin{requirement}[No ad hoc exponentiation]
\label{req:no-ad-hoc-exponentiation}
The replacement
\begin{equation}
 1+\cK^{(1)}\longmapsto \exp\!\bigl[\cK^{(1)}\bigr]
\label{eq:forbidden-eikonal-exponent}
\end{equation}
is prohibited unless the exponentiated object, color algebra, subtraction
scheme, and accuracy are specified and the result is benchmarked against the
explicit second order.  Non-Abelian commutators, crossed exchanges, and
state-changing insertions may not be hidden in a scalar exponent.
\end{requirement}

\subsection{Order-by-order unitarity and cuts}

Let \(S=1+iT\) in the retained regulated space.  Exact unitarity gives
\begin{equation}
 -i(T-T^\dagger)=T^\dagger T.
\label{eq:unitarity-operator}
\end{equation}
At order \(g^n\), the discontinuity of the virtual amplitude is related to
all retained cuts whose lower-order amplitudes multiply to \(g^n\).  Define
\begin{equation}
 \cU_n
 =-i\bigl(T^{(n)}-T^{(n)\dagger}\bigr)
 -\sum_{m=1}^{n-1}T^{(m)\dagger}T^{(n-m)}.
\label{eq:unitarity-residual-order-n}
\end{equation}
A complete truncated channel has \(\cU_n=0\) up to regulator, omitted-sector,
and numerical errors.

The eikonal Wilson line is not an isolated scattering matrix, so
\cref{eq:unitarity-residual-order-n} is applied to the enlarged amplitude
containing the struck-parton operator, spectator state, and retained soft
radiation.  It is nevertheless a powerful test: an imaginary term with no
corresponding retained cut signals an incomplete intermediate-state basis or
an inconsistent pole prescription.

\subsection{Convergence observables}

The eikonal expansion is accepted through order \(N\) only if stability is
shown for matrix elements and projected observables.  For an observable
\(\cO_A\), define
\begin{align}
 \delta^{\rm ord}_{A,N}
 & =
 \frac{
  \left|\cO_A[U_{[N]}]-\cO_A[U_{[N-1]}]\right|
 }{
  \max\{\left|\cO_A[U_{[N]}]\right|,\cO_{A,\rm ref}\}
 },
\label{eq:eikonal-order-residual}\\
 \delta^{\rm rep}_{A,N}
 & =
 \frac{
  \left|\cO_A[U_{[N]}^{\rm D}]
       -\cO_A[U_{[N]}^{\rm M}]_{\text{expanded to }N}\right|
 }{
  \max\{\left|\cO_A\right|,\cO_{A,\rm ref}\}
 }.
\label{eq:representation-residual}
\end{align}
These are reported separately from basis, regulator, rapidity, and numerical
errors.  A stable Sivers moment does not establish stability of a different
Boer--Mulders or gluon color projection.

\section{Process maps, factorization status, and Glauber structure}
\label{sec:process-factorization}

\subsection{A process object is more than a staple sign}

A process declaration is the typed object
\begin{equation}
\begin{aligned}
 \mathfrak p=\bigl(&
 \text{name},\text{external representations},\text{measured final state},
 \text{power accuracy},\text{perturbative order},\\
 &\gamma_q,[\gamma_g,\gamma'_g],
 C_f^{\mathfrak p},C_d^{\mathfrak p},
 \mathfrak f_{\rm status},\mathfrak g_{\rm Glauber}\bigr).
\end{aligned}
\label{eq:process-object}
\end{equation}
The link data are outputs of the factorization analysis, not labels inferred
from the words ``SIDIS,'' ``Drell--Yan,'' or ``gluon process.''  A hard
channel may also project a linear combination of independent gluon color
structures,
\begin{equation}
 W_{g}^{\mathfrak p}
 =C_f^{\mathfrak p}W_g^{(f)}
  +C_d^{\mathfrak p}W_g^{(d)},
\label{eq:process-fd-combination}
\end{equation}
where the weights belong to \(\Proc\), not to the universal parent registry.

\subsection{Factorization-status lattice}

We use the following fail-closed status set:
\begin{longtable}{L{0.19\textwidth}L{0.72\textwidth}}
\toprule
Status & Meaning and permitted use\\
\midrule
\endhead
\texttt{ESTABLISHED} &
A leading-power theorem with declared observable, regulator/scheme, link
structure, soft subtraction, and demonstrated cancellation or factorization
of Glauber contributions.  Production predictions are permitted within the
stated accuracy.\\
\texttt{CONDITIONAL} &
A derivation exists under named assumptions, restricted color channels,
kinematics, or measurement inclusiveness.  Predictions must carry those
conditions.\\
\texttt{EXPLORATORY} &
A process-motivated link or color map is used as a model study without a
complete factorization theorem.  Results may be benchmarks or sensitivities,
not universal extractions.\\
\texttt{BROKEN} &
Known color entanglement or Glauber exchange prevents the assumed product of
independent universal TMDs.  The unsupported composition is rejected.\\
\texttt{UNASSESSED} &
No sufficient analysis has been registered.  No process-level prediction is
permitted.\\
\bottomrule
\end{longtable}
The statuses are partially ordered by evidential content, not by numerical
precision.  A successful fit cannot promote \texttt{EXPLORATORY} to
\texttt{ESTABLISHED}.

\subsection{Reference color-singlet processes}

For sufficiently inclusive color-singlet SIDIS and Drell--Yan observables in
the TMD region, the future/past quark-staple relation and its Sivers sign
change follow from the corresponding factorization analyses
\cite{Collins2002,JiMaYuan2004,JiMaYuanSpin2004}.  They provide the primary
process-reversal benchmark:
\begin{align}
 \gamma_q^{\rm SIDIS}&=\gamma_+,
 &\gamma_q^{\rm DY}&=\gamma_-,
\label{eq:sidis-dy-process-map}\\
 f_{1T}^{\perp q,\rm SIDIS}
 &=-f_{1T}^{\perp q,\rm DY},
 &h_1^{\perp q,\rm SIDIS}
 &=-h_1^{\perp q,\rm DY},
\label{eq:sidis-dy-two-signs}
\end{align}
within common operator, Fourier, and spin conventions.  The theorem is not a
statement that every hadronic process with an initial- or final-state colored
parton admits the same two links.

\subsection{Glauber region and cancellation contract}

A Glauber momentum has parametrically small longitudinal products compared
with its transverse virtuality.  In a representative two-direction problem,
\begin{equation}
 \ell^+\ell^-\ll \ell_T^2,
\label{eq:glauber-region}
\end{equation}
and ordinary soft contour deformations may be pinched.  A valid process proof
must show one of the following:
\begin{enumerate}
\item the Glauber region cancels after the declared sum over cuts;
\item it is absorbed into the same Wilson/soft objects used in the theorem;
\item it is represented by explicit Glauber operators whose matrix elements
      factorize in the asserted manner;
\item the proposed TMD factorization fails.
\end{enumerate}
SCET formulations with explicit Glauber operators make this obligation
particularly transparent \cite{RothsteinStewart2016}.

Define a process-specific residual
\begin{equation}
 \delta_{\rm G}^{\mathfrak p}
 =\cA_{\rm full}^{\mathfrak p}
 -\cA_{\rm coll}^{\mathfrak p}
 -\cA_{\soft}^{\mathfrak p}
 -\cA_{\rm G,assigned}^{\mathfrak p}
 +\cA_{\rm overlap}^{\mathfrak p}.
\label{eq:glauber-residual}
\end{equation}
The statement \(\delta_{\rm G}^{\mathfrak p}=0\) requires a proof or a
controlled calculation; it is not a configuration flag.

\subsection{Known failure pattern}

Hadron--hadron production of nearly back-to-back colored final states can
contain color entanglement that is not representable by a product of two
independent universal TMD correlators.  Explicit counterexamples show that
process-dependent gauge links alone do not always repair the factorization
\cite{CollinsQiu2007,RogersMulders2010}.  The process registry therefore
rejects a schematic expression
\begin{equation}
 \dd\sigma
 \stackrel{\rm unsupported}{=}
 H\otimes \Phi_{A}^{[\gamma_A]}
   \otimes \Phi_{B}^{[\gamma_B]}
\label{eq:unsupported-factorization}
\end{equation}
when the color flow requires one entangled matrix element.

\begin{requirement}[Factorization-status propagation]
\label{req:factorization-status-propagation}
Every observable, plot, likelihood term, and exported table must retain
\(\mathfrak f_{\rm status}\), the supporting theorem or analysis identifier,
the link/color map, and the stated power/perturbative accuracy.  A downstream
routine may not drop this metadata when converting a correlator into a number.
\end{requirement}

\subsection{Gluon-sensitive channels}

Gluon TMD processes require an ordered link pair and process-specific color
weights.  The registry must answer, for each candidate channel:
\begin{enumerate}
\item Which field-strength color indices are connected to which eikonals?
\item Which of \([+,+],[-,-],[+,-],[-,+]\), or a more complicated link
      topology is generated?
\item Which \(f\)- and \(d\)-type contractions enter the hard color trace?
\item Does the measurement preserve or trace over the color information needed
      for factorization?
\item What Glauber cancellation or obstruction is known?
\end{enumerate}
No default ``gluon Sivers link'' is defined in the absence of these answers.

\section{Gauge-link dependence of orbital motion and final-state torque}
\label{sec:oam-torque}

\subsection{Straight and staple definitions}

A Wigner/GTMD definition of partonic orbital angular momentum retains its
Wilson path.  At a matched scale and scheme, define
\begin{equation}
 L_z^a[\gamma]
 =\int \dd x\,\dd^2\kT\,\dd^2\bD\,
 \crossT{\bD}{\kT}\,
 \rho_{LU}^{a[\gamma]}(x,\kT,\bD).
\label{eq:path-dependent-oam}
\end{equation}
A straight connector and a staple connector need not yield the same operator:
\begin{equation}
 \Delta L_z^a
 =L_z^a[\gamma_{\staple}]-L_z^a[\gamma_{\straight}].
\label{eq:oam-difference}
\end{equation}
In common interpretations, the straight-link quantity is related to kinetic
OAM and the staple-link quantity to canonical OAM in a specified gauge-link
and boundary prescription; the precise identification is scheme and operator
convention dependent \cite{LorcePasquini2011,Hatta2012,Burkardt2013}.

\subsection{Path derivative and field-strength insertion}

An infinitesimal deformation of a Wilson path produces a field-strength
insertion.  Schematically, for a deformation \(\delta\gamma\),
\begin{align}
 \delta U_\gamma[b,a]
 ={}&-ig\,A_\mu(b)\delta b^\mu U_\gamma
 +ig\,U_\gamma A_\mu(a)\delta a^\mu
\notag\\
 &-ig\int_0^1\dd s\,
 U_{\gamma[b,\gamma(s)]}
 F_{\mu\nu}(\gamma(s))
 U_{\gamma[\gamma(s),a]}
 \dot\gamma^\nu(s)\,\delta\gamma^\mu(s).
\label{eq:path-derivative-field-strength}
\end{align}
The bulk term is the operator origin of the accumulated color force or torque
between the straight and staple paths.  It also shows why path dependence is
dynamical when \(F_{\mu\nu}\ne0\): homotopy data alone do not determine
\cref{eq:oam-difference}.

\subsection{Shared OAM blocks and naive-\texorpdfstring{\(T\)}{T}-odd projections}

The same \(L_z\)-resolved Fock amplitudes enter OAM moments and T-odd
interference, but their kernels differ.  We may write
\begin{align}
 L_z^a[\gamma]
 &=\sum_{\alpha\beta}
   \psi_\alpha^*\,
   \cL_{\alpha\beta}^{a[\gamma]}\,
   \psi_\beta,
\label{eq:oam-kernel}\\
 f_{1T}^{\perp a[\gamma]}
 &=\sum_{\alpha\beta}
   \psi_\alpha^*\,
   \cS_{\alpha\beta}^{a[\gamma]}\,
   \psi_\beta.
\label{eq:sivers-kernel}
\end{align}
A common pair \((\alpha,\beta)\) can contribute to both, producing correlated
sensitivity to microscopic OAM content.  It does not imply
\(\cL\propto\cS\).

\begin{principle}[No universal lensing identity]
\label{pr:no-universal-lensing}
A relation between a Sivers function and a transverse distortion from a GPD
requires a model-dependent factorization of spatial distortion and final-state
interaction, often called chromodynamic lensing.  Volume~III may test such a
relation within a specified model, but it does not promote
\begin{equation}
 f_{1T}^{\perp q}\overset{?}{\propto} E^q
\label{eq:forbidden-universal-lensing}
\end{equation}
to an operator identity.
\end{principle}

\section{Sum rules, positivity domains, and global consistency}
\label{sec:sum-rules}

\subsection{Angular and collinear limits}

A naive-\(T\)-odd leading-twist term carries an explicit transverse harmonic.
For example,
\begin{equation}
 \Phi^{[\gamma^+]}(x,\kT;S_T)
 \supset
 -\frac{\epsilon_T^{ij}k_T^iS_T^j}{M}
 f_{1T}^{\perp}(x,k_T^2).
\label{eq:sivers-harmonic}
\end{equation}
Consequently the unweighted azimuthal integral of the full modulation
vanishes.  This does not imply that the scalar coefficient
\(f_{1T}^{\perp}\) or its weighted moment vanishes.  The renormalized first
moment is a separate, generally ultraviolet-sensitive operator quantity.

At the matched collinear level, no ordinary twist-two PDF carries the Sivers
or Boer--Mulders harmonic.  Their transverse moments instead match to
quark--gluon or tri-gluon twist-three correlators.  The adapter is defined in
\cref{app:twist3}.

\subsection{Net transverse-momentum sum rule}

Transverse momentum conservation implies a Burkardt-type sum rule for the
average transverse momentum generated by the Sivers mechanism.  In a common
moment convention,
\begin{equation}
 \sum_{a=q,\bar q,g}
 \int_0^1\dd x\,
 f_{1T}^{\perp(1)a}(x;\mu)
 =0,
\label{eq:burkardt-sum-rule}
\end{equation}
up to the declared sign, normalization, renormalization, and endpoint
conventions \cite{Burkardt2004,GoekeMeissnerMetzSchlegel2006}.  The sum is a
constraint on the complete state.  A valence-only calculation is not expected
to satisfy it if the omitted sea and gluon sectors carry compensating
transverse momentum.

Define the closure residual
\begin{equation}
 \delta_{\rm BSR}
 =\sum_{a=q,\bar q,g}
 \int_{x_{\min}}^{x_{\max}}\dd x\,
 f_{1T}^{\perp(1)a}(x)
 +\delta_{x\to0}+\delta_{x\to1}
 +\delta_{\rm match}.
\label{eq:burkardt-residual}
\end{equation}
Endpoint and matching terms are reported rather than silently set to zero.

\subsection{Color neutrality and Ward closure}

For a physical color-singlet state,
\begin{equation}
 Q^a_{\rm color}|N\rangle=0.
\label{eq:color-neutrality-repeated}
\end{equation}
The sum of eikonal attachments to all retained charged lines must satisfy the
corresponding Ward identity.  Let \(\cK_i^{\mu a}\) be the attachment to line
\(i\).  Then
\begin{equation}
 \ell_\mu\sum_i \cK_i^{\mu a}
 =\delta_{\rm Ward}^a,
\label{eq:ward-residual-sum}
\end{equation}
where \(\delta_{\rm Ward}^a\to0\) as constrained-field, transverse-link,
sector, and numerical completeness are restored.  A nonzero result may not be
renormalized away by changing the T-odd normalization.

\subsection{Positivity is layer dependent}

For a fixed physical link class, a complete forward cut correlator can admit
a Gram representation
\begin{equation}
 \Phi_{\alpha\beta}^{[\gamma]}
 =\sum_X A_{\alpha X}^{[\gamma]}
       A_{\beta X}^{[\gamma]*},
\label{eq:link-specific-gram}
\end{equation}
and is positive semidefinite.  Three qualifications are essential:
\begin{enumerate}
\item positivity applies to the complete joint parton--target matrix, not to
      every scalar function independently;
\item an off-forward GTMD or Wigner quasidistribution is not pointwise
      positive;
\item a fixed-order expansion of \cref{eq:link-specific-gram}, or a
      soft-subtracted scheme object, need not be positive term by term.
\end{enumerate}
The implementation never clips a negative eigenvalue.  It records whether the
object is incomplete, perturbatively truncated, or genuinely inconsistent.

For two retained amplitude channels \(A\) and \(B\), Cauchy--Schwarz gives
\begin{equation}
 \left|\operatorname{Im}\sum_X A_XB_X^*\right|
 \le
 \left(\sum_X|A_X|^2\right)^{1/2}
 \left(\sum_X|B_X|^2\right)^{1/2}.
\label{eq:phase-cauchy-schwarz}
\end{equation}
This supplies a structural bound on an interference-generated T-odd
projection without imposing a separate phenomenological ceiling on each
named function.

\subsection{Common-state correlation obligation}

Let \(\theta\) denote the Hamiltonian and Wilson-sector parameters.  A
variation \(\delta\theta_a\) must propagate to every projection containing
the corresponding vertex:
\begin{equation}
 J_{Aa}=\frac{\partial \cO_A}{\partial\theta_a}.
\label{eq:phase-jacobian}
\end{equation}
In particular, a parameter controlling an eikonal infrared scale, a quark--
gluon vertex, or an OAM mixing block affects Sivers, Boer--Mulders, twist-three
moments, path-dependent OAM, and any gluon channels to which the same
interaction contributes.  Assigning independent copies of that parameter to
each TMD would destroy the principal predictive content of this volume.

\section{Computational realization and formal interfaces}
\label{sec:software}

\subsection{Core objects}

The computational representation follows the physical decomposition rather
than hiding all information in one complex array.
\begin{longtable}{L{0.27\textwidth}L{0.64\textwidth}}
\toprule
Object & Responsibility\\
\midrule
\endhead
\texttt{WilsonPathId} &
C1 path identity: ordered segments, endpoints, orientation, finite or
semi-infinite extent, transverse closure, representation, and inverse.\\
\texttt{PolePrescription} &
The sign and variable of every eikonal or light-front energy denominator;
principal-value and delta pieces; regulator deformation; provenance of the
\(i0\).\\
\texttt{EikonalVertex} &
Momentum, Lorentz, color, representation, path-segment, helicity, and sector
map for one Wilson attachment.\\
\texttt{IntermediateStateId} &
Fock sector, basis member, momentum routing, conserved quantum numbers,
energy, regulator level, and completeness weight of a retained state.\\
\texttt{CutSource} &
A cut denominator or eikonal pole, the cut set to which it belongs, and the
rule preventing duplicate imaginary support.\\
\texttt{RescatteringKernel} &
Sparse \(\Amp\)-layer map assembled from vertices, resolvents, color words,
cut data, and Wilson order.  It acts on Fock amplitudes rather than named
TMDs.\\
\texttt{WilsonExpansion} &
Dyson or Magnus representation, strict/unexpanded order, path-ordering
metadata, convergence residuals, and conversion adapters.\\
\texttt{ColorContractionId} &
Fundamental/adjoint representation and, for three-adjoint gluon poles,
independent \(f\)- and \(d\)-type contractions with normalization.\\
\texttt{AbsorptiveAmplitude} &
Real and imaginary components with explicit cut provenance, source/target
states, link orientation, and Hermitian/time-reversal partners.\\
\shortstack[l]{\texttt{RapiditySubtraction}\\\texttt{Contract}} &
Rapidity regulator, zero-bin/overlap prescription, soft factor, square-root
allocation, \((\mu,\zeta)\), and the boundary-only or joint-soft route.\\
\texttt{PhaseBudget} &
Microscopic Wilson, constrained-field, transverse-link, soft, matching, and
excluded duplicate contributions with a closure residual.\\
\texttt{ProcessLinkMap} &
Process definition, hard color flow, quark link, ordered gluon link pair,
\(f/d\) weights, power accuracy, factorization status, and Glauber evidence.\\
\shortstack[l]{\texttt{EikonalConvergence}\\\texttt{Manifest}} &
Wilson order, basis, Fock, regulator, rapidity, quadrature, and representation
residuals for every reported matrix element and projection.\\
\bottomrule
\end{longtable}

\subsection{Typed data flow}

The intended computation is
\begin{align}
 |\Psi_N\rangle
 &\xrightarrow{\ \cO_{a}^{[\Gamma]}\ }
 W_{a}^{(0)}
 \xrightarrow{\ \cK_{\eta,[N]}\ }
 W_{a,\unsub}^{[\gamma]}
 \xrightarrow{\ \cS^{-1/2}\cR_{\rap}\cZ_{\UV}\ }
 W_{a,\ren}^{[\gamma]},
\notag\\
 W_{a,\ren}^{[\gamma]}
 &\xrightarrow{\ \Red_A\ }
 F_A^{[\gamma]}
 \xrightarrow{\ \Proc_{\mathfrak p}\ }
 \cO_{\mathfrak p}.
\label{eq:software-data-flow}
\end{align}
The first two arrows are \(\Amp\)-layer maps.  Soft subtraction and
regulator/scheme conversion are \(\Match\)-layer maps.  Projection is a
\(\Red\)-layer map, and the final process assembly is a \(\Proc\)-layer map.
No implicit conversion between these layers is permitted.

The link-odd projection is formed only after both orientations exist:
\begin{equation}
 W_{\odd}
 =\frac12\left(
   W_{\ren}^{[\gamma_+]}
  -\Theta^{-1}W_{\ren}^{[\gamma_-]}\Theta
 \right),
\label{eq:software-link-odd}
\end{equation}
where \(\Theta\) is the full antiunitary convention adapter.  Subtracting two
arrays with path labels \texttt{FUTURE} and \texttt{PAST} but without the
spin/momentum adapter is a type error.

\subsection{Sparse block structure}

A kernel block is indexed by
\begin{equation}
 \mathfrak b=
 (\nu_{\rm in},\nu_{\rm out},
  J^z_{\rm in},J^z_{\rm out},
  L^z_{\rm in},L^z_{\rm out},
  R_{c,\rm in},R_{c,\rm out},
  \Gamma,\gamma,\mathfrak c,N,r).
\label{eq:kernel-block-id}
\end{equation}
Forbidden blocks are absent by symmetry.  A nonzero block must carry the
intertwiner that proves color, flavor, and angular compatibility.  The sparse
representation allows all contributing \(L_z\) and Fock interferences to be
retained without forming one dense matrix across the full basis.

\subsection{Principal-value and cut quadrature}

A numerical integral containing
\begin{equation}
 \frac{F(\ell)}{x(\ell)+i0\sigma}
\label{eq:numerical-singular-integral}
\end{equation}
is evaluated as two declared operations,
\begin{align}
 I_{\PV}
 &=\PV\int\dd\ell\,\frac{F(\ell)}{x(\ell)},
\label{eq:numerical-pv}\\
 I_{\Disc}
 &=-i\sigma\pi
   \int_{x(\ell)=0}
   \frac{\dd\Sigma(\ell)}{|\nabla x(\ell)|}
   F(\ell).
\label{eq:numerical-cut-surface}
\end{align}
The cut-surface Jacobian is explicit.  A finite imaginary regulator used for
numerical smoothing must converge to \cref{eq:numerical-cut-surface}; it may
not define the physical width of the phase.

Recommended checks are:
\begin{enumerate}
\item symmetric-domain principal-value cancellation for odd test functions;
\item convergence against analytic Hilbert transforms;
\item stability as the finite regulator tends to zero;
\item equality of explicit delta-surface and contour-deformation evaluations;
\item grid refinement near pinched or endpoint support.
\end{enumerate}

\subsection{Automatic differentiation and parameter sharing}

The kernel may be differentiated with respect to physical parameters,
\begin{equation}
 \frac{\partial W_A}{\partial\theta_a}
 =\inner{\partial_a\Psi}{\cO_A U_\gamma\Psi}
 +\inner{\Psi}{(\partial_a\cO_A)U_\gamma\Psi}
 +\inner{\Psi}{\cO_A(\partial_a U_\gamma)\Psi}
 +\inner{\Psi}{\cO_AU_\gamma\partial_a\Psi},
\label{eq:full-phase-derivative}
\end{equation}
with adjoint methods used when the state is an implicit Hamiltonian
eigenvector.  Derivatives must include the parameter dependence of the state,
not only of the explicit rescattering kernel.  A ``frozen wave function''
derivative is a named approximation.

\subsection{Serialization and reproducibility}

Every exported kernel or projected table stores:
\begin{enumerate}
\item Volume~I state and regulator identity;
\item Volume~II operator and recoil identity;
\item path, pole, color, Wilson order, and cut identities;
\item rapidity and soft-subtraction contract;
\item process/factorization status if a process map has been applied;
\item parameter/member identity and all truncation coordinates;
\item source-code commit, configuration hash, and numerical tolerances.
\end{enumerate}
A complex array without this manifest is not an authoritative Volume~III
object.

\section{Formal requirements and property-based tests}
\label{sec:tests}

\subsection{Stable requirement identifiers}

The following identifiers are normative and should appear in the software
coverage manifest.
\begin{longtable}{L{0.14\textwidth}L{0.77\textwidth}}
\toprule
ID & Requirement\\
\midrule
\endhead
\texttt{V3.PATH.1} & Bare transport composes on matching endpoint fibers and
satisfies inverse/adjoint identities.\\
\texttt{V3.PATH.2} & Staple orientation, finite/infinite extent, transverse
closure, and representation are explicit identity fields.\\
\texttt{V3.POLE.1} & Every eikonal denominator has a process-derived \(i0\)
sign and a tested distributional decomposition.\\
\texttt{V3.POLE.2} & Every light-front energy denominator has an independent
pole identity; equal algebraic denominators are not deduplicated without cut
provenance.\\
\texttt{V3.CUT.1} & The imaginary part is reproducible from an explicit
intermediate-state or eikonal cut.\\
\texttt{V3.CUT.2} & A cut ledger prevents the same on-shell support from
entering both the Wilson and resolvent terms twice.\\
\texttt{V3.WARD.1} & Longitudinal polarization replacement and total color
attachment satisfy the declared Ward test.\\
\texttt{V3.TIME.1} & The antiunitary adapter maps future to past operators,
including spin, momenta, endpoints, representations, and complex conjugation.\\
\texttt{V3.TIME.2} & Link-even and link-odd projections have the correct
future/past behavior after renormalization.\\
\texttt{V3.ZERO.1} & Setting the gauge coupling or rescattering kernel to zero
removes every naive-\(T\)-odd projection.\\
\texttt{V3.QUARK.1} & Sivers and Boer--Mulders use distinct operator and
helicity projectors on a shared kernel.\\
\texttt{V3.QUARK.2} & Quark, flavor, and positive-\(x\) antiquark identities
remain explicit through the phase calculation.\\
\texttt{V3.GLUON.1} & Gluon operators retain an ordered pair of adjoint links.\\
\texttt{V3.GLUON.2} & \(f\)- and \(d\)-type color tensors are independent,
orthogonal channels and are not inferred from a TMD name.\\
\texttt{V3.SOFT.1} & The microscopic phase budget closes after soft/rapidity
overlap subtraction.\\
\texttt{V3.SOFT.2} & Boundary-only and joint-microscopic-soft routes are
exclusive unless an explicit conversion/subtraction is supplied.\\
\texttt{V3.ORD.1} & One-gluon results reproduce an analytic benchmark.\\
\texttt{V3.ORD.2} & Higher orders report Dyson/Magnus, strict/unexpanded, and
observable-level convergence residuals.\\
\texttt{V3.PROC.1} & Every process carries a theorem status, link/color map,
power accuracy, and Glauber assessment.\\
\texttt{V3.PROC.2} & Known color-entangled compositions fail closed.\\
\texttt{V3.SUM.1} & The common state reports the transverse-momentum sum-rule
residual with endpoint and matching terms.\\
\texttt{V3.OAM.1} & Straight/staple OAM differences retain path identity and
are not converted into a universal lensing relation.\\
\texttt{V3.REPRO.1} & Every authoritative result has a complete manifest and
can be regenerated from a clean environment.\\
\bottomrule
\end{longtable}

\subsection{Algebraic tests}

For random admissible finite paths and color matrices, property-based tests
must verify
\begin{align}
 U_{\gamma_3\circ(\gamma_2\circ\gamma_1)}
 &=U_{(\gamma_3\circ\gamma_2)\circ\gamma_1},
\label{eq:path-associativity-test}\\
 U_{\gamma^{-1}}U_\gamma&=\bbone,
\label{eq:path-inverse-test}\\
 U_\gamma^\dagger U_\gamma&=\bbone
\quad\text{for a Hermitian connection and exact transport},
\label{eq:path-unitarity-test}\\
 f^{abc}d^{abc}&=0.
\label{eq:fd-test}
\end{align}
Finite-order Wilson approximants are tested against the expected order of the
unitarity defect rather than required to be exactly unitary.

\subsection{Pole and cut tests}

For smooth compactly supported test functions \(\varphi\), verify numerically
\begin{equation}
 \lim_{\epsilon\to0^+}
 \int\dd x\,
 \frac{\varphi(x)}{x+i\sigma\epsilon}
 =\PV\int\dd x\,\frac{\varphi(x)}{x}
 -i\sigma\pi\varphi(0).
\label{eq:distribution-numerical-test}
\end{equation}
The future/past pair must share the principal-value part and have opposite
discontinuities.  A deliberately reversed \(i0\) sign must fail the
link-reversal test.

For the LF resolvent, a finite discrete model with known eigenvalues
\(E_X\) must satisfy
\begin{equation}
 \operatorname{Im}
 \sum_X\frac{N_X}{E_i-E_X+i0}
 =-\pi\sum_X N_X\delta(E_i-E_X)
\label{eq:discrete-resolvent-test}
\end{equation}
after replacing the delta distribution by the declared finite-volume rule.

\subsection{Symmetry and zero tests}

The mandatory injected failures are:
\begin{enumerate}
\item replace the future pole by the past pole while keeping the future path;
\item reverse endpoints without conjugating the color word;
\item apply complex conjugation without the spin/momentum time-reversal map;
\item remove one required \(\Delta L_z=1\) amplitude from a Sivers block;
\item use the vector projector for Boer--Mulders or the chiral-odd projector
      for Sivers;
\item exchange the order of a mixed gluon link pair;
\item relabel an \(f\)-type color contraction as \(d\)-type;
\item include a microscopic soft ladder and an external CS kernel without
      overlap subtraction;
\item apply a TMD product formula to a process registered as
      \texttt{BROKEN};
\item count one on-shell pole in both the eikonal and intermediate-state cut.
\end{enumerate}
Each injection must be rejected before an authoritative numerical output is
written.

\subsection{One-gluon and higher-order tests}

The one-gluon calculation must pass:
\begin{enumerate}
\item analytic eikonal-pole sign and normalization;
\item Abelian color limit;
\item SU(3) fundamental and adjoint color factors;
\item zero-coupling and zero-cut limits;
\item future/past sign reversal;
\item Ward closure after all declared attachments are included;
\item agreement of direct complex integration and PV-plus-cut evaluation.
\end{enumerate}
At second order, the strict Dyson result must agree with the expanded Magnus
result through \(g^2\), and the connected color commutator must be nonzero for
a non-Abelian benchmark.

\subsection{Soft and scheme tests}

For two rapidity regulators or subtraction schemes connected by a known
finite map, verify
\begin{equation}
 H_{\mathfrak s}W_{\mathfrak s}^{[\gamma]}
 -H_{\mathfrak s'}W_{\mathfrak s'}^{[\gamma]}
 =\cO(\text{first omitted perturbative order}).
\label{eq:scheme-invariance-test}
\end{equation}
The same conversion must act on future and past links.  Their ratio or
link-odd difference may not acquire a spurious finite sign from inconsistent
soft-factor square roots.

The rapidity consistency residual is
\begin{equation}
 \delta_{\mu\zeta}
 =
 \frac{\partial}{\partial\ln\mu}
 \frac{\partial}{\partial\ln\sqrt\zeta}
 \ln W
 -
 \frac{\partial}{\partial\ln\sqrt\zeta}
 \frac{\partial}{\partial\ln\mu}
 \ln W.
\label{eq:rapidity-integrability-residual}
\end{equation}
It must vanish to the stated order after matching.

\section{Analytic benchmark suite}
\label{sec:benchmarks}

\subsection{Benchmark A: a semi-infinite Abelian line}

Take a commuting gauge field with one Fourier mode,
\begin{equation}
 A_\mu(x)=\epsilon_\mu e^{-i\ell\cdot x}.
\label{eq:abelian-mode}
\end{equation}
For \(\gamma_\eta(\lambda)=a+\eta\lambda v\),
\begin{equation}
 U_{\eta}^{(1)}
 =-g\,e^{-i\ell\cdot a}
 \frac{v\cdot\epsilon}{v\cdot\ell-i0\eta}.
\label{eq:abelian-line-benchmark}
\end{equation}
This benchmark fixes the path, Fourier, coupling, and pole signs before any
hadron model is introduced.  Its future/past difference is
\begin{equation}
 U_+^{(1)}-U_-^{(1)}
 =-2i\pi g\,e^{-i\ell\cdot a}
 (v\cdot\epsilon)\delta(v\cdot\ell).
\label{eq:abelian-line-difference}
\end{equation}

\subsection{Benchmark B: two-state light-front cut}

Let \(|0\rangle,|1\rangle\) have light-front energies \(E_0,E_1\), and let
\(V_{10}=v\), \(O_{01}=o\).  The rescattering term is
\begin{equation}
 \cA_\sigma
 =\frac{o v}{E_0-E_1+i0\sigma}.
\label{eq:two-state-amplitude}
\end{equation}
Then
\begin{equation}
 \operatorname{Im}\cA_\sigma
 =-\sigma\pi o v\,\delta(E_0-E_1)
\label{eq:two-state-cut}
\end{equation}
for real \(ov\), while an off-shell discrete system has zero absorptive part
until a continuum or finite-volume cut prescription is introduced.  This
prevents a numerical eigensolver broadening from being mistaken for physical
rescattering.

\subsection{Benchmark C: minimal OAM interference}

Use two real radial amplitudes with azimuthal weights \(L_z=0\) and \(L_z=1\),
\begin{equation}
 \psi(\kT)=u_0(k_T)+e^{i\phi_k}u_1(k_T).
\label{eq:minimal-oam-state}
\end{equation}
Let the link-odd kernel contribute \(i\eta K_{01}\) between the two blocks.
The odd interference is
\begin{equation}
 \cI_{01}^{[\eta]}
 =2\eta\,u_0u_1\,K_{01}\,\sin\phi_k.
\label{eq:minimal-oam-interference}
\end{equation}
It vanishes when \(u_0\), \(u_1\), or \(K_{01}\) is disabled and reverses
under \(\eta\to-\eta\).  Replacing the target-spin projector by the active-
quark transverse-spin projector gives a distinct observable with the same
radial state; this is the minimal no-proportionality benchmark for Sivers and
Boer--Mulders.

\subsection{Benchmark D: SU(3) color orthogonality}

With Gell-Mann normalization
\(\Tr(t^at^b)=\tfrac12\delta^{ab}\), construct random adjoint tensors and
project them with \(if^{abc}\) and \(d^{abc}\).  The cross term must vanish,
\begin{equation}
 (if^{abc})^*d^{abc}=0,
\label{eq:fd-orthogonality-benchmark}
\end{equation}
while the norms agree with \cref{eq:f-norm,eq:d-norm} below.  Swapping a mixed
link pair must not automatically swap the color projector.

\subsection{Benchmark E: subtraction overlap}

Choose a regulated scalar model in which a soft region \(S(\delta)\) is
present both in the unsubtracted collinear graph and in a separate soft graph:
\begin{align}
 W_{\unsub}(\delta)&=W_{\rm finite}+S(\delta),
\notag\\
 \frac12\ln \cS(\delta)&=S(\delta)+\cO(g^4).
\label{eq:scalar-soft-benchmark}
\end{align}
Then the order-\(g^2\) subtracted boundary is regulator independent,
\begin{equation}
 W_{\subtr}=W_{\rm finite}+\cO(g^4).
\label{eq:scalar-soft-cancellation}
\end{equation}
Adding the same \(S(\delta)\) again through an external evolution kernel is an
injected double-counting failure.

\subsection{Benchmark F: process rejection}

Construct a toy hard color tensor that connects the color indices of two
incoming correlators in one trace.  It cannot be represented as
\(C_A\Phi_A\,C_B\Phi_B\).  The process registry must return
\texttt{BROKEN} or \texttt{UNASSESSED}, and the standard two-TMD process map
must refuse evaluation.  This benchmark tests architecture, not a numerical
smallness assumption.

\section{Contract with Volume IV: matched spin-1 nuclear dynamics}
\label{sec:volume4-contract}

\subsection{Nucleon output consumed by the nuclear theory}

For every nucleon species \(N=p,n\), Volume~III exports
\begin{equation}
 \mathfrak W_{a/N}^{(3)}=
 \left(
 W_{a/N,\even}^{[\gamma]},
 W_{a/N,\odd}^{[\gamma]},
 \mathfrak o_a,
 \mathfrak t_{\rm Wilson},
 \mathfrak B,
 \mathfrak p_{\rm allowed},
 \operatorname{Cov}_{\rm micro}
 \right),
\label{eq:volume3-export}
\end{equation}
where \(\mathfrak B\) is the phase budget and
\(\mathfrak p_{\rm allowed}\) is the set of process maps whose factorization
status permits use.  Proton and neutron flavor, state member, Wilson member,
and regulator identities remain resolved.

Volume~IV may insert \cref{eq:volume3-export} into the spin-resolved nuclear
matrix element, but it may not replace the path/color data by a generic
``T-odd nucleon function.''

\subsection{Partonic staple versus nuclear rescattering}

The partonic Wilson line associated with the hard scattering and the
rescattering of a struck nucleon through the nuclear environment are not
automatically the same object.  We distinguish
\begin{align}
 U_{\gamma}^{\rm parton}
 &: \text{color transport in the partonic operator},
\notag\\
 \cU_{\rm nuc}
 &: \text{coherent propagation among nuclear Fock components}.
\label{eq:parton-nuclear-transport}
\end{align}
Their scales and degrees of freedom may overlap after matching.  Volume~IV
must therefore provide a subtraction or factorization map
\begin{equation}
 \cZ_{\rm parton\leftrightarrow nuc}
\label{eq:parton-nuclear-matching}
\end{equation}
whenever the same soft exchange can appear in both descriptions.

\begin{requirement}[No nuclear phase multiplication]
\label{req:no-nuclear-phase-multiplication}
A tensor-polarized deuteron T-odd distribution may not be constructed by
multiplying an unpolarized nucleon TMD by an independently fitted nuclear
phase.  Tensor, vector, and unpolarized channels must descend from the
spin-resolved nuclear amplitude acting on the link-resolved nucleon operator.
\end{requirement}

\subsection{Coherence remains at amplitude level}

Let \(\alpha,\beta\) label \(NN\), \(NN\pi\), \(\Delta\Delta\), \(6q\), or
other retained nuclear sectors.  The link-dependent deuteron parent is
\begin{equation}
 W_{a/D}^{[\gamma]}
 =\sum_{\alpha,\beta}
 \inner{\Psi_D^\alpha}
 {\widehat\cO_{a,\alpha\beta}^{[\gamma]}
  \Psi_D^\beta}.
\label{eq:volume4-full-amplitude}
\end{equation}
Off-diagonal terms remain until the observable no longer resolves their
coherence.  Only then may a partial trace generate a completely positive
reduced map.  Coherent shadowing or exchange currents may not be inserted as
an incoherent CP map merely to preserve positivity.

\subsection{Nuclear link-reversal test}

If the nuclear interaction is invariant under the relevant antiunitary
operation and the process factorization produces only the standard staple
reversal, then
\begin{equation}
 W_{a/D,\odd}^{[\gamma_+]}
 =-W_{a/D,\odd}^{[\gamma_-]}.
\label{eq:nuclear-link-reversal}
\end{equation}
A failure of \cref{eq:nuclear-link-reversal} must be traced to one of:
\begin{enumerate}
\item a convention adapter error;
\item an incomplete nuclear state or current;
\item an additional process link or color structure;
\item a factorization-breaking/Glauber contribution;
\item numerical or regulator inconsistency.
\end{enumerate}
It must not be repaired by sign-flipping the final named TMD.

\section{Contract with Volume V: matching, evolution, and observables}
\label{sec:volume5-contract}

\subsection{Matching datum}

Volume~V receives, at a declared low-resolution boundary,
\begin{equation}
 \Bigl(W_{a/h,\subtr}^{[\gamma]},
  \cR_{\UV},\cR_{\rap},\cS_{\soft},
  \mu_0,\zeta_0,
  r,N,\mathfrak B
 \Bigr)
\label{eq:volume5-input}
\end{equation}
and constructs
\begin{equation}
 W_{a/h,\TMD}^{[\gamma]}(\mu,\zeta)
 =\cU_a(\mu,\zeta;\mu_0,\zeta_0)
  \otimes
  \cZ_{\LF\to\TMD}^{a[\gamma]}
  \otimes
  W_{a/h,\subtr}^{[\gamma]}.
\label{eq:volume5-matching-evolution}
\end{equation}
The matching coefficient may depend on the parton representation, operator,
transverse rank, and path class.  The universal leading-power Collins--Soper
kernel must not erase the link/color identity of the boundary.

\subsection{Evolution preserves link reversal}

Suppose future and past operators share the same renormalization and rapidity
kernel.  Then
\begin{equation}
 W_{\odd}^{[+]}(\mu_0,\zeta_0)
 =-W_{\odd}^{[-]}(\mu_0,\zeta_0)
\quad\Longrightarrow\quad
 W_{\odd}^{[+]}(\mu,\zeta)
 =-W_{\odd}^{[-]}(\mu,\zeta).
\label{eq:evolution-preserves-sign}
\end{equation}
A violation after evolution is therefore a scheme, implementation, or
factorization inconsistency unless an additional process-dependent operator
has entered.

\subsection{Twist-three overlap region}

At perturbative transverse momentum, the TMD and collinear twist-three
representations describe an overlap region.  Volume~V must compare
\begin{equation}
 W_{\TMD}^{\text{expanded}}
 \quad\text{with}\quad
 W_{\rm twist3}^{\text{collinear}}
\label{eq:tmd-twist3-overlap}
\end{equation}
using the convention adapter in \cref{app:twist3}.  The equality is at the
level of a matched observable or coefficient expansion, not a bare equality
between a cutoff moment and a renormalized twist-three function
\cite{JiQiuVogelsangYuan2006,JiQiuVogelsangYuanSIDIS2006}.

\subsection{Process kernels and the \texorpdfstring{\(W+Y\)}{W+Y} construction}

A process prediction is
\begin{equation}
 \dd\sigma_{\mathfrak p}
 =W_{\mathfrak p}
  +Y_{\mathfrak p},
 \qquad
 Y_{\mathfrak p}
 =\dd\sigma_{\rm FO}^{\mathfrak p}
  -\left[W_{\mathfrak p}\right]_{\rm expanded\ to\ FO},
\label{eq:volume5-wy}
\end{equation}
only when \(\mathfrak f_{\rm status}\) permits the factorized
\(W_{\mathfrak p}\).  A high-\(q_T\) discrepancy may not be absorbed by
retuning the microscopic rescattering phase.

\subsection{No-double-counting square}

The complete handoff must make the following square commute to the stated
accuracy:
\begin{equation}
\begin{tikzcd}[column sep=large,row sep=large]
 W_{\micro,\unsub}^{[\gamma]}
 \arrow[r,"\text{soft/rapidity subtraction}"]
 \arrow[d,"\text{strict expansion}"']
 & W_{\micro,\subtr}^{[\gamma]}
 \arrow[d,"\cZ_{\LF\to\TMD}"]
 \\
 W_{\rm FO,\micro}^{[\gamma]}
 \arrow[r,"\text{overlap subtraction}"']
 & W_{\TMD}^{[\gamma]}\big|_{\rm expanded}
\end{tikzcd}
\label{eq:no-double-counting-square}
\end{equation}
Failure identifies a missing matching term, regulator mismatch, or duplicate
soft region.  It is not an uncertainty band.

\section{Staged implementation program}
\label{sec:stages}

\subsection{Relation to the current repository migration}

C2 migrates the accepted phenomenological boundary to native typed reductions
and an enforceable provenance graph.  It does not introduce the microscopic
Hamiltonian or the Wilson dynamics of this volume.  Volume~III therefore
supplies the normative specification for a later C3 package after the actual
C2 report fixes repository-level APIs and remaining gaps.

The recommended computational stages are:
\begin{longtable}{L{0.14\textwidth}L{0.76\textwidth}}
\toprule
Package & Scientific and computational content\\
\midrule
\endhead
\texttt{C3A} &
Native eikonal path segments, pole prescriptions, distribution identities,
rapidity-regulator identities, and analytic Abelian benchmarks.\\
\texttt{C3B} &
Light-front resolvent, intermediate-state and cut objects, one-gluon
\texttt{RescatteringKernel}, Ward tests, and cut-provenance ledger on a small
controlled Fock basis.\\
\texttt{C3C} &
Quark link-even/link-odd GTMD projections, Sivers and Boer--Mulders projector
separation, flavor/antiquark propagation, and twist-three moment adapters.\\
\texttt{C3D} &
Ordered gluon link pairs, adjoint Wilson vertices, \(f/d\) color projections,
small-\(x\) specializations, and independent color-channel tests.\\
\texttt{C3E} &
Soft/rapidity subtraction contract, phase budget, process/factorization
registry, Glauber status, and rejection of unsupported process compositions.\\
\texttt{C3F} &
Second-order Dyson/Magnus benchmark, observable-level eikonal convergence,
common-state covariance, sum-rule closure, complete reports, and a frozen
interface for Volumes IV--V.\\
\bottomrule
\end{longtable}

\subsection{Pilot hierarchy}

The first pilot should use the smallest state that can test the physics:
\begin{equation}
 \Hil_{\rm pilot}
 =\Hil_{\rm active+spectator}
  \oplus\Hil_{\rm active+spectator+g},
\label{eq:pilot-hilbert-space}
\end{equation}
with at least two OAM blocks.  It must reproduce the one-gluon benchmark,
future/past reversal, a nonzero Sivers-like interference, and a distinct
Boer--Mulders-like projection.  Only after these pass should the engine be
connected to the full Volume~I nucleon state.

For gluons, the first pilot must contain enough color structure to distinguish
\(f\) and \(d\), not merely an Abelian adjoint label.  A calculation that
produces one generic ``gluon T-odd'' array does not satisfy C3D.

\subsection{Calibration and prediction split}

A useful first inference program may calibrate a restricted common
rescattering sector to selected quark Sivers information, while withholding:
\begin{itemize}
\item Boer--Mulders data or moments;
\item flavor combinations not directly fitted;
\item one or both gluon color sectors;
\item path-dependent OAM or torque observables;
\item tensor-polarized nuclear projections in Volume~IV.
\end{itemize}
The withheld quantities count as predictions only when no independent TMD-
specific normalization or width is introduced after calibration.

\section{Acceptance criteria and scientific status}
\label{sec:acceptance}

\subsection{Completion criteria for the Volume III pilot}

The dynamical Wilson-line pilot is complete only when all of the following
hold:
\begin{enumerate}
\item The bare path, endpoint fibers, representation, staple orientation,
      transverse closure, and inverse are explicit and tested.
\item Every eikonal and light-front energy denominator has a derived pole
      prescription and cut provenance.
\item The one-gluon phase is reproduced from an analytic benchmark and an
      explicit intermediate-state cut.
\item Ward identities close after all required spectator, constrained-field,
      and transverse-link attachments are included.
\item Setting the gauge coupling, eikonal kernel, or required OAM block to zero
      removes the corresponding naive-\(T\)-odd projection.
\item Future/past reversal is exact after the full antiunitary adapter, before
      and after soft subtraction and evolution.
\item Sivers and Boer--Mulders are different operator projections of a common
      rescattering kernel and are not linked by an imposed proportionality.
\item Positive-\(x\) antiquark and all quark flavor identities survive the
      calculation.
\item Gluon correlators retain ordered links and independent \(f\)- and
      \(d\)-type color contractions.
\item The phase budget closes and the microscopic/soft overlap is subtracted
      exactly once.
\item Dyson/Magnus and eikonal-order convergence are reported for each claimed
      observable family.
\item The process registry rejects a known color-entangled factorization
      pattern and propagates status into every allowed prediction.
\item The transverse-momentum sum-rule residual, endpoint estimate, and
      matching residual are reported from a common state.
\item All regulator, basis, Fock, Wilson-order, rapidity, and numerical
      uncertainties remain separate axes.
\item At least one observable family not used to calibrate the common phase is
      successfully predicted.
\end{enumerate}

\subsection{What completion does and does not establish}

Passing these criteria establishes a controlled dynamical rescattering model
within the declared regulated Hamiltonian, Fock basis, Wilson order, and TMD
matching route.  It does not establish:
\begin{itemize}
\item an all-orders factorization theorem for every process;
\item exact continuum QCD;
\item convergence of the later spin-1 nuclear expansion;
\item a universal relation among all T-odd functions;
\item a universal gluon \(f/d\) mixture;
\item the absence of higher-twist or Glauber corrections outside the declared
      process and kinematics.
\end{itemize}

\subsection{Recommended formal foundation}

The adopted Volume~III object is
\begin{quote}
\emph{a gauge-covariant, path- and process-resolved operator-valued
rescattering kernel acting on the common regulated light-front state, with
explicit eikonal and light-front cuts, independent quark/gluon color and spin
projections, matched soft/rapidity subtraction, controlled Wilson-order
convergence, and fail-closed factorization status.}
\end{quote}
Its value is not that Wilson-line terminology creates a phase.  Its value is
that the phase must be generated by a specified gauge interaction and cut,
propagate coherently across all operator projections, and survive every
matching and process step without losing its origin.

\appendix

\section{Distribution identities and sign audit}
\label{app:distributions}

\subsection{Sokhotski--Plemelj convention}

For \(\varphi\in C_c^\infty(\mathbb R)\),
\begin{equation}
 \lim_{\epsilon\to0^+}
 \int_{-\infty}^{\infty}\dd x\,
 \frac{\varphi(x)}{x\pm i\epsilon}
 =\PV\int_{-\infty}^{\infty}\dd x\,
 \frac{\varphi(x)}{x}
 \mp i\pi\varphi(0).
\label{eq:sokhotski-plemelj}
\end{equation}
Thus
\begin{equation}
 \frac1{x-i0\eta}
 =\PV\frac1x+i\eta\pi\delta(x).
\label{eq:sokhotski-eta}
\end{equation}
This fixes \cref{eq:eikonal-distribution-identity}.

\subsection{Semi-infinite Fourier integral}

For \(\epsilon>0\),
\begin{align}
 I_\eta(x;\epsilon)
 &=\int_0^\infty\dd\lambda\,
   \eta e^{-i\eta\lambda x-\epsilon\lambda}
 =\frac{\eta}{\epsilon+i\eta x}
 =-\frac{i}{x-i\eta\epsilon}.
\label{eq:semi-infinite-regulated}
\end{align}
Taking \(\epsilon\to0^+\) yields
\begin{equation}
 I_\eta(x)
 =-i\PV\frac1x+\eta\pi\delta(x).
\label{eq:semi-infinite-pv-delta}
\end{equation}
Multiplication by the Wilson exponent \(-ig\) gives a convention-dependent
real/imaginary assignment.  The invariant audit is the sign change of the
discontinuity under \(\eta\to-\eta\), after all common factors are included.

\subsection{Momentum-flow conversion}

If the gluon momentum is redefined \(\ell\to-\ell\), then
\begin{equation}
 \frac{1}{v\cdot\ell-i0\eta}
 \longrightarrow
 -\frac{1}{v\cdot\ell+i0\eta}.
\label{eq:momentum-flow-pole}
\end{equation}
The extra minus sign combines with the vertex and momentum-routing
conventions.  The software may convert between routings only through an
adapter that transforms the numerator, color ordering, and cut identity as
well as the denominator.

\section{SU(\texorpdfstring{\(N_c\)}{Nc}) color dictionary}
\label{app:color}

Use
\begin{equation}
 \Tr(t^at^b)=T_F\delta^{ab},
 \qquad T_F=\frac12,
\label{eq:generator-normalization}
\end{equation}
with
\begin{equation}
 [t^a,t^b]=if^{abc}t^c,
 \qquad
 \{t^a,t^b\}=\frac1{N_c}\delta^{ab}\bbone+d^{abc}t^c.
\label{eq:fd-definitions}
\end{equation}
The standard contractions are
\begin{align}
 f^{amn}f^{bmn}&=N_c\delta^{ab},
\label{eq:f-contraction}\\
 d^{amn}d^{bmn}&=\frac{N_c^2-4}{N_c}\delta^{ab},
\label{eq:d-contraction}\\
 f^{amn}d^{bmn}&=0.
\label{eq:fd-contraction-orthogonality}
\end{align}
Consequently,
\begin{align}
 f^{abc}f^{abc}&=N_c(N_c^2-1),
\label{eq:f-norm}\\
 d^{abc}d^{abc}&=\frac{(N_c^2-4)(N_c^2-1)}{N_c}.
\label{eq:d-norm}
\end{align}
For \(N_c=3\), these norms are \(24\) and \(40/3\), respectively.
Normalized projectors can be defined as
\begin{align}
 (P_f)^{abc}
 &=\frac{if^{abc}}{\sqrt{N_c(N_c^2-1)}},
\notag\\
 (P_d)^{abc}
 &=\sqrt{\frac{N_c}{(N_c^2-4)(N_c^2-1)}}\,d^{abc}.
\label{eq:normalized-fd-projectors}
\end{align}
The factor of \(i\) is part of the convention for the antisymmetric color
channel and must be included in Hermiticity and time-reversal adapters.

\section{Light-front cut ledger}
\label{app:cuts}

For a term with denominators \(D_j=\Delta E_j+i0\sigma_j\), define a cut
record
\begin{equation}
\begin{aligned}
 \mathfrak c=\bigl(&
 \{j_1,\ldots,j_m\},\{\Delta E_{j_r}=0\},\{\sigma_{j_r}\},
 \text{intermediate states},\\
 &\text{phase-space Jacobian},\text{source graph/path word}\bigr).
\end{aligned}
\label{eq:cut-record}
\end{equation}
The discontinuity generated by cutting one denominator is
\begin{equation}
 \Disc_{D_j}
 \prod_{r}\frac1{D_r}
 =-2\pi i\sigma_j\delta(\Delta E_j)
  \prod_{r\ne j}\frac1{\Delta E_r+i0\sigma_r}.
\label{eq:single-cut-product}
\end{equation}
Multiple cuts are generated by repeated discontinuities with the appropriate
combinatorics and ordering.  The ledger canonicalizes the physical cut set
without erasing graph or path provenance.  A cut represented by a soft
factor and by an explicit microscopic graph is marked as an overlap relation,
not as two additive contributions.

The following fields are mandatory:
\begin{longtable}{L{0.28\textwidth}L{0.62\textwidth}}
\toprule
Field & Meaning\\
\midrule
\endhead
\texttt{cut\_source} & \texttt{EIKONAL\_POLE},
\texttt{LF\_ENERGY}, \texttt{CONSTRAINED\_FIELD},
\texttt{SOFT\_OVERLAP}, or another registered source.\\
\texttt{support\_equations} & Equations defining the on-shell or pole
surface.\\
\texttt{state\_ids} & Complete intermediate Fock/basis states on the cut.\\
\texttt{jacobian} & Measure induced on the cut surface.\\
\texttt{orientation} & Future/past and momentum-flow convention.\\
\texttt{color\_word} & Ordered generators or invariant tensor.\\
\texttt{overlap\_group} & Identifier joining alternative representations of
the same physical support.\\
\texttt{status} & Included, subtracted, benchmark-only, or omitted with an
error estimate.\\
\bottomrule
\end{longtable}

\section{Twist-three moment adapter}
\label{app:twist3}

\subsection{Quark Sivers moment}

Define the rank-one transverse moment
\begin{equation}
 f_{1T}^{\perp(1)q[\gamma]}(x;\mu)
 =\int\dd^2\kT\,
 \frac{k_T^2}{2M^2}
 f_{1T}^{\perp q[\gamma]}(x,k_T^2;\mu),
\label{eq:sivers-first-moment}
\end{equation}
with the understanding that the integral requires ultraviolet
renormalization or a regulated \(b\)-space derivative.  A convention adapter
defines the diagonal Qiu--Sterman function by
\begin{equation}
 T_F^q(x,x;\mu)
 =\sigma_{\rm QS}
  2M\,f_{1T}^{\perp(1)q[\gamma_+]}(x;\mu),
\qquad
 \sigma_{\rm QS}\in\{+1,-1\},
\label{eq:qiu-sterman-adapter}
\end{equation}
where \(\sigma_{\rm QS}\) stores the Fourier, epsilon-tensor, Wilson-link,
and correlator conventions.  It is never inferred from a dataset name.

The TMD and twist-three descriptions overlap in the perturbative tail, but
\cref{eq:qiu-sterman-adapter} is not applied to an unrenormalized cutoff
moment.  The matching order and subtraction scheme are part of the adapter.

\subsection{Boer--Mulders and chiral-odd twist three}

The Boer--Mulders first moment similarly matches a chiral-odd quark--gluon
correlator,
\begin{equation}
 h_1^{\perp(1)q[\gamma]}
 \xleftrightarrow{\ \cZ_{\rm BM}^{\rm twist3}\ }
 T_F^{(\sigma)q},
\label{eq:boer-mulders-twist3}
\end{equation}
with a different spin projector and matching kernel from
\cref{eq:qiu-sterman-adapter}.  The two twist-three functions are not related
by the common eikonal interaction alone.

\subsection{Tri-gluon channels}

The two gluonic-pole color contractions define independent tri-gluon
correlators,
\begin{equation}
 T_G^{(f)}(x,x;\mu),
 \qquad
 T_G^{(d)}(x,x;\mu),
\label{eq:tri-gluon-functions}
\end{equation}
with separate convention and normalization adapters.  A process may select
\(C_f^{\mathfrak p}T_G^{(f)}+C_d^{\mathfrak p}T_G^{(d)}\), but the parent
store never replaces the pair by one universal function.

\section{Machine-readable schemas}
\label{app:schemas}

\subsection{Rescattering-kernel record}

A minimal JSON-like schema is
\begin{verbatim}
{
  "schema": "deuteron_wigner.volume3.rescattering_kernel.v1",
  "source_operator_id": "...",
  "source_sector": "...",
  "target_sector": "...",
  "wilson_path_id": "...",
  "pole_prescriptions": ["..."],
  "intermediate_state_ids": ["..."],
  "cut_ids": ["..."],
  "color_contraction_id": "...",
  "wilson_order": 1,
  "series_representation": "DYSON_STRICT",
  "uv_regulator": "...",
  "rapidity_regulator": "...",
  "soft_contract_id": "...",
  "parameter_member": "...",
  "provenance_node": "..."
}
\end{verbatim}

\subsection{Process-link record}

\begin{verbatim}
{
  "schema": "deuteron_wigner.volume3.process_link_map.v1",
  "process_id": "...",
  "observable_id": "...",
  "power_accuracy": "LEADING_POWER",
  "perturbative_order": "...",
  "quark_path_id": "...",
  "gluon_link_pair_id": "...",
  "f_weight": "...",
  "d_weight": "...",
  "factorization_status": "ESTABLISHED|CONDITIONAL|EXPLORATORY|BROKEN|UNASSESSED",
  "factorization_reference": "...",
  "glauber_status": "...",
  "allowed": true
}
\end{verbatim}

\subsection{Phase-budget record}

\begin{verbatim}
{
  "schema": "deuteron_wigner.volume3.phase_budget.v1",
  "operator_id": "...",
  "microscopic_wilson_terms": ["..."],
  "intermediate_state_terms": ["..."],
  "constrained_field_terms": ["..."],
  "transverse_link_terms": ["..."],
  "soft_overlap_terms": ["..."],
  "matching_terms": ["..."],
  "excluded_duplicates": ["..."],
  "closure_residual": 0.0,
  "tolerance": 0.0,
  "status": "PASS|FAIL"
}
\end{verbatim}

\section{Compact acceptance checklist}
\label{app:checklist}

Before any Volume~III output is used in a fit or nuclear calculation, answer:
\begin{enumerate}
\item Is the path geometry complete, including transverse closure?
\item Is every pole sign derived from the same Fourier and momentum-flow
      convention?
\item Can every imaginary contribution be traced to a cut?
\item Are Ward and color-singlet tests closed at the retained truncation?
\item Does zero rescattering remove the claimed T-odd projection?
\item Does the full antiunitary adapter give exact future/past reversal?
\item Are Sivers and Boer--Mulders separate projectors?
\item Are quark flavors and antiquarks explicit?
\item Are gluon link order and \(f/d\) color independent?
\item Is the soft/rapidity overlap subtracted exactly once?
\item Are Wilson-order and regulator convergence documented?
\item Does the process have an admissible factorization status?
\item Are twist-three moments matched rather than equated to bare cutoff
      integrals?
\item Are sum-rule and phase-budget residuals reported?
\item Is the complete provenance manifest attached?
\end{enumerate}
A negative answer blocks authoritative use and identifies the next formal or
computational task.

\small
\begin{thebibliography}{99}

\bibitem{BrodskyHwangSchmidt2002}
S.~J.~Brodsky, D.~S.~Hwang, and I.~Schmidt,
``Final-state interactions and single-spin asymmetries in semi-inclusive deep
inelastic scattering,''
Phys.\ Lett.\ B \textbf{530}, 99 (2002), arXiv:hep-ph/0201296.

\bibitem{Collins2002}
J.~C.~Collins,
``Leading-twist single-transverse-spin asymmetries: Drell--Yan and deep
inelastic scattering,''
Phys.\ Lett.\ B \textbf{536}, 43 (2002), arXiv:hep-ph/0204004.

\bibitem{BelitskyJiYuan2003}
A.~V.~Belitsky, X.~Ji, and F.~Yuan,
``Final-state interactions and gauge-invariant parton distributions,''
Nucl.\ Phys.\ B \textbf{656}, 165 (2003), arXiv:hep-ph/0208038.

\bibitem{BoerMuldersPijlman2003}
D.~Boer, P.~J.~Mulders, and F.~Pijlman,
``Universality of T-odd effects in single spin and azimuthal asymmetries,''
Nucl.\ Phys.\ B \textbf{667}, 201 (2003), arXiv:hep-ph/0303034.

\bibitem{BomhofMuldersPijlman2004}
C.~J.~Bomhof, P.~J.~Mulders, and F.~Pijlman,
``Gauge link structure in quark--quark correlators in hard processes,''
Phys.\ Lett.\ B \textbf{596}, 277 (2004), arXiv:hep-ph/0406099.

\bibitem{JiMaYuan2004}
X.~Ji, J.-P.~Ma, and F.~Yuan,
``QCD factorization for semi-inclusive deep-inelastic scattering at low
transverse momentum,''
Phys.\ Rev.\ D \textbf{71}, 034005 (2005), arXiv:hep-ph/0404183.

\bibitem{JiMaYuanSpin2004}
X.~Ji, J.-P.~Ma, and F.~Yuan,
``QCD factorization for spin-dependent cross sections in DIS and Drell--Yan
processes at low transverse momentum,''
Phys.\ Lett.\ B \textbf{597}, 299--308 (2004), arXiv:hep-ph/0405085.

\bibitem{JiQiuVogelsangYuan2006}
X.~Ji, J.-W.~Qiu, W.~Vogelsang, and F.~Yuan,
``A unified picture for single transverse-spin asymmetries in hard
processes,''
Phys.\ Rev.\ Lett.\ \textbf{97}, 082002 (2006), arXiv:hep-ph/0602239.

\bibitem{JiQiuVogelsangYuanSIDIS2006}
X.~Ji, J.-W.~Qiu, W.~Vogelsang, and F.~Yuan,
``Single transverse-spin asymmetry in semi-inclusive deep inelastic
scattering,''
Phys.\ Lett.\ B \textbf{638}, 178--186 (2006), arXiv:hep-ph/0604128.

\bibitem{Meissner2009}
S.~Mei{\ss}ner, A.~Metz, and M.~Schlegel,
``Generalized parton correlation functions for a spin-\(\tfrac12\) hadron,''
JHEP \textbf{08}, 056 (2009), arXiv:0906.5323.

\bibitem{EchevarriaGTMD2016}
M.~G.~Echevarria, A.~Idilbi, K.~Kanazawa, C.~Lorc\'e, A.~Metz,
B.~Pasquini, and M.~Schlegel,
``Proper definition and evolution of generalized transverse momentum
distributions,''
Phys.\ Lett.\ B \textbf{759}, 336--341 (2016), arXiv:1602.06953.

\bibitem{EchevarriaSoft2016}
M.~G.~Echevarria, I.~Scimemi, and A.~Vladimirov,
``Universal transverse momentum dependent soft function at NNLO,''
Phys.\ Rev.\ D \textbf{93}, 054004 (2016), arXiv:1511.05590.

\bibitem{Collins2011}
J.~C.~Collins,
\emph{Foundations of Perturbative QCD},
Cambridge University Press (2011).

\bibitem{BuffingMulders2011}
M.~G.~A.~Buffing and P.~J.~Mulders,
``Gauge links for transverse-momentum-dependent correlators at tree level,''
JHEP \textbf{07}, 065 (2011), arXiv:1105.4804.

\bibitem{BuffingGluon2013}
M.~G.~A.~Buffing, A.~Mukherjee, and P.~J.~Mulders,
``Generalized universality of definite rank gluon transverse momentum
dependent correlators,''
Phys.\ Rev.\ D \textbf{88}, 054027 (2013), arXiv:1306.5897.

\bibitem{Boer2016}
D.~Boer, S.~Cotogno, T.~van Daal, P.~J.~Mulders, A.~Signori, and Y.-J.~Zhou,
``Gluon and Wilson loop TMDs for hadrons of spin \(\leq1\),''
JHEP \textbf{10}, 013 (2016), arXiv:1607.01654.

\bibitem{BoerEchevarriaMuldersZhou2016}
D.~Boer, M.~G.~Echevarria, P.~J.~Mulders, and J.~Zhou,
``Single spin asymmetries from a single Wilson loop,''
Phys.\ Rev.\ Lett.\ \textbf{116}, 122001 (2016), arXiv:1511.03485.

\bibitem{CollinsQiu2007}
J.~C.~Collins and J.-W.~Qiu,
``\(k_T\) factorization is violated in production of high-transverse-momentum
particles in hadron--hadron collisions,''
Phys.\ Rev.\ D \textbf{75}, 114014 (2007), arXiv:0705.2141.

\bibitem{RogersMulders2010}
T.~C.~Rogers and P.~J.~Mulders,
``No generalized TMD-factorization in the hadro-production of high transverse
momentum hadrons,''
Phys.\ Rev.\ D \textbf{81}, 094006 (2010), arXiv:1001.2977.

\bibitem{RothsteinStewart2016}
I.~Z.~Rothstein and I.~W.~Stewart,
``An effective field theory for forward scattering and factorization
violation,''
JHEP \textbf{08}, 025 (2016), arXiv:1601.04695.

\bibitem{BuffingMuldersColor2014}
M.~G.~A.~Buffing and P.~J.~Mulders,
``Color entanglement for azimuthal asymmetries in the Drell--Yan process,''
Phys.\ Rev.\ Lett.\ \textbf{112}, 092002 (2014), arXiv:1309.4681.

\bibitem{LorcePasquini2011}
C.~Lorc\'e and B.~Pasquini,
``Quark Wigner distributions and orbital angular momentum,''
Phys.\ Rev.\ D \textbf{84}, 014015 (2011), arXiv:1106.0139.

\bibitem{Hatta2012}
Y.~Hatta,
``Notes on the orbital angular momentum of quarks in the nucleon,''
Phys.\ Lett.\ B \textbf{708}, 186 (2012), arXiv:1111.3547.

\bibitem{Burkardt2013}
M.~Burkardt,
``Parton orbital angular momentum and final state interactions,''
Phys.\ Rev.\ D \textbf{88}, 014014 (2013), arXiv:1205.2916.

\bibitem{Burkardt2004}
M.~Burkardt,
``Sivers mechanism for gluons,''
Phys.\ Rev.\ D \textbf{69}, 091501 (2004), arXiv:hep-ph/0402014.

\bibitem{GoekeMeissnerMetzSchlegel2006}
K.~Goeke, S.~Mei{\ss}ner, A.~Metz, and M.~Schlegel,
``Checking the Burkardt sum rule for the Sivers function by model
calculations,''
Phys.\ Lett.\ B \textbf{637}, 241 (2006), arXiv:hep-ph/0601133.

\bibitem{Magnus1954}
W.~Magnus,
``On the exponential solution of differential equations for a linear
operator,''
Commun.\ Pure Appl.\ Math.\ \textbf{7}, 649 (1954).

\bibitem{Gatheral1983}
J.~G.~M.~Gatheral,
``Exponentiation of eikonal cross sections in non-Abelian gauge theories,''
Phys.\ Lett.\ B \textbf{133}, 90 (1983).

\bibitem{FrenkelTaylor1984}
J.~Frenkel and J.~C.~Taylor,
``Non-Abelian eikonal exponentiation,''
Nucl.\ Phys.\ B \textbf{246}, 231 (1984).

\end{thebibliography}

\end{document}
